
\medterm{OAB} Alternate index label; see \textit{Overactive bladder}.

\medterm{Oat Cell Carcinoma} Alternate historical term; see \textit{Small-Cell Carcinoma}.

\medterm{OB/GYN} Abbreviation for obstetrics and gynecology, the medical specialties that care for pregnancy, childbirth, the period after birth, and the reproductive system.



\medterm{Obesity} A chronic disease involving abnormal or excessive accumulation of body fat that may impair health. It is linked with a higher risk of conditions such as type 2 diabetes, high blood pressure, heart disease, stroke, sleep apnea, osteoarthritis, and some cancers. Many factors can contribute, including genes, metabolism, hormones, medicines, eating patterns, activity, environment, and social conditions; it is not caused by one factor or by lack of willpower alone. Body mass index (BMI) is commonly used to classify obesity in adults, but it does not measure body fat directly and is interpreted with other information.

\medterm{Obesity Hypoventilation Syndrome} A breathing disorder in which excess body weight contributes to too little breathing and a buildup of carbon dioxide while awake. It often occurs with sleep apnea and may cause sleepiness, headaches, breathlessness, low oxygen, and heart strain.

\medterm{Obex Region} A small area at the lower end of the brain’s fourth fluid-filled space, where it narrows into the central canal of the spinal cord. It lies near centers controlling breathing, swallowing, and vomiting.

\medterm{Objective Response Rate} The percentage of people in a cancer study whose tumors shrink by a predefined amount or disappear during treatment. It does not show how long the response lasts or whether treatment improves survival.

\medterm{Objects Or Insects In Ear} An object or insect in the ear canal can cause pain, hearing loss, bleeding, discharge, or a sensation of movement. Do not insert tools or liquids; removal by a clinician is safest when the object is deep, sharp, chemical, or near the eardrum.

\synonyms
Ear Object In (Objects Or Insects In Ear)


\medterm{Oblique Fracture} A break in a bone that runs diagonally across it.

\medterm{Oblique Plane} A slice or view through the body or an organ at an angle other than straight across. Medical imaging can use oblique planes to show structures that are difficult to see in standard horizontal, vertical, or front-to-back views.

\medterm{Obliteration} Complete closing, destruction, or disappearance of a space, passage, or hollow channel in the body. It may result from scarring, inflammation, growth, injury, or treatment and can prevent normal movement of air, blood, fluid, or other contents.
\medterm{Observation} Planned monitoring of a person, symptom, disease, or test result without immediate treatment or a procedure. Follow-up timing and signs that require reassessment should be defined.

\medterm{Obsession} A recurrent, intrusive, unwanted thought, image, or urge that causes distress and is difficult to control.
\medterm{Obsessions} Repeated, unwanted thoughts, images, or urges that enter the mind despite attempts to stop them and cause distress. Common examples involve contamination, harm, mistakes, symmetry, or taboo subjects; having an obsession does not mean a person will act on it.

\medterm{Obsessive Compulsive Disorder} Alternate index label for Obsessive-Compulsive Disorder.

\medterm{Obsessive-Compulsive And Related Disorders} A group of mental disorders involving recurrent intrusive thoughts, repetitive behaviors or mental acts, preoccupation with perceived appearance flaws, difficulty discarding possessions, or repetitive body-focused behaviors. Examples include obsessive-compulsive disorder, body dysmorphic disorder, hoarding disorder, trichotillomania, and excoriation disorder.

\medterm{Obsessive-Compulsive Disorder} A mental-health disorder involving unwanted, repeated thoughts, images, or urges and repetitive behaviors or mental acts performed to reduce distress or prevent a feared event. The symptoms are time-consuming or cause clinically significant distress or impairment.

\medterm{Obsessive-Compulsive Personality Disorder} A long-lasting pattern of excessive concern with order, rules, perfection, and control. It can make a person inflexible, slow to complete tasks, or unable to delegate and can interfere with relationships or work.

\medterm{Obstetric Anal Sphincter Injury} A severe tear during childbirth that extends into the muscles controlling the anus. It may cause pain, difficulty controlling gas or stool, or later incontinence and needs prompt examination, repair, and follow-up.
\medterm{Obstetric Cholestasis} Alternate index label; see \textit{Cholestasis of pregnancy}.

\medterm{Obstetric Fistula} A hole that forms between the birth canal and the bladder or rectum after prolonged, obstructed labour without timely medical care. It causes continuous leakage of urine or stool and may lead to skin infections, kidney problems, and social isolation. Most cases can be prevented with timely maternity care, and many can be repaired with surgery.

\medterm{Obstetric Nursing} Obstetric nursing provides care during pregnancy, labor, birth, and the postpartum period, including monitoring, education, support, and recognition of complications.

\medterm{Obstetric Ultrasound} An ultrasound scan during pregnancy used to check the baby's age, growth, anatomy, position, heartbeat, and sometimes the placenta and fluid.

\medterm{Obstetrical Forceps} A surgical instrument used in selected vaginal births to guide the baby’s head when delivery needs assistance. Safe use requires the correct position, an experienced clinician, and assessment of risks and alternatives.

\medterm{Obstetrician} A medical doctor who cares for pregnant people during pregnancy, labor, birth, and the postpartum period. Obstetricians manage routine and high-risk pregnancies, perform deliveries, and treat complications such as high blood pressure, diabetes, bleeding, or preterm labor.

\medterm{Obstetrics} The branch of medicine concerned with care during pregnancy, labor, birth, and the weeks after birth, including routine care and treatment of complications affecting the pregnant person or baby.

\medterm{Obstipation} Severe constipation with inability to pass stool or gas. It can occur with a blockage in the intestine and may cause abdominal swelling, pain, vomiting, or serious illness requiring urgent assessment.
\medterm{Obstructed Infected Urinary System} A blocked urinary tract with an infection behind the blockage. It is a medical emergency because infection can spread quickly and the blockage often must be relieved urgently.

\medterm{Obstructed Labour} Failure of the fetus to pass through the birth canal despite adequate uterine contractions, commonly because of disproportion, abnormal presentation, or an obstructing lesion.

\medterm{Obstruction} A blockage that prevents normal movement through a body passage, such as an airway, intestine, bile duct, blood vessel, or urinary tract. Causes include swelling, stones, scars, tumors, clots, and foreign objects.

\medterm{Obstructive Azoospermia} A cause of male infertility in which no sperm enters the semen because a reproductive tube is blocked or missing, although the testes make sperm normally. Causes include a missing vas deferens, vasectomy, infection-related scarring, blocked ejaculatory ducts, or surgical injury. Testicle size and hormone levels are often normal. Treatment may reopen the tube or collect sperm for use with in-vitro fertilization.

\synonyms
OA; Excretory Azoospermia

\medterm{Obstructive Hydrocephalus} A buildup of cerebrospinal fluid caused by a blockage in its normal flow through the brain's ventricles.

\medterm{Obstructive Jaundice} Yellowing of the skin or eyes caused by blockage of bile flow from the liver. It may produce dark urine, pale stools, itching, abdominal pain, or fever and can result from a stone, narrowing, inflammation, or tumor.

\medterm{Obstructive Nephropathy} Kidney damage caused by a blockage that prevents urine from draining normally. Stones, tumors, scarring, an enlarged prostate, or birth-related problems can cause it; untreated blockage may lead to infection, permanent loss of kidney function, or kidney failure.

\medterm{Obstructive Parotitis} Inflammation of a saliva gland near the ear because its drainage tube is blocked, often by a stone. Painful swelling usually worsens during meals, and trapped saliva can lead to repeated infection and tenderness.

\medterm{Obstructive Shock} A life-threatening drop in blood flow caused by a physical blockage, such as a large blood clot in the lungs, pressure around the heart, or a collapsed lung under pressure.

\medterm{Obstructive Sleep Apnea} Repeated narrowing or closing of the upper airway during sleep, causing pauses in breathing, snoring, fragmented sleep, daytime tiredness, and sometimes reduced oxygen levels.

\synonyms
Sleep Apnea, Obstructive

\medterm{Obstructive Sleep Apnea - Adults} Obstructive sleep apnea is repeated upper-airway collapse during sleep, causing pauses in breathing, oxygen fluctuation, snoring, fragmented sleep, and daytime impairment.


\medterm{Obstructive Uropathy} A blockage anywhere between the kidneys and the urethra that slows or stops urine flow. Stones, scarring, an enlarged prostate, tumors, or pressure from nearby tissue may cause it and can damage the kidneys if not relieved.

\medterm{Obtunded} Having reduced alertness and slowed responses but still responding to strong voice or touch. It can result from medicines, illness, metabolic problems, poisoning, or brain disease and needs assessment.

\medterm{Obturator} Relating to a structure that closes an opening. In the pelvis, the obturator nerve helps control muscles that draw the thigh inward and carries feeling from part of the inner thigh; the obturator artery supplies nearby muscles.

\medterm{Obturator Hernia} A rare hernia in which part of the intestine pushes through an opening in the pelvic bone. It can block the bowel and may cause pain in the inner thigh.

\medterm{Obtuse Marginal Artery} A branch of the circumflex coronary artery that runs along the outer border of the heart and supplies part of the left ventricle. Its size and number vary between people; blockage can cause a heart attack.


\medterm{Occipital} Relating to the back of the head, the occipital bone, or the occipital lobe. These structures protect and support the brain and help process visual information and head or neck movement.
\medterm{Occipital Bone} The bone at the back and base of the skull. It surrounds the opening where the spinal cord joins the brain and connects with the first neck vertebra.

\medterm{Occipital Diastasis} Widening or separation of a joint between bones at the back of a newborn’s skull. It may follow difficult birth or abnormal development, and imaging assesses its extent and associated injury.

\medterm{Occipital Hematoma} A collection of blood in the scalp or tissues over the back of the head, often after trauma; a rapidly enlarging swelling or associated neurologic symptoms requires urgent assessment.

\medterm{Occipital Lobe} The part of the brain at the back of the head that receives and interprets visual signals. Injury or disease may cause missing areas of vision, difficulty recognizing objects, or visual hallucinations.

\medterm{Occipital Neuralgia} Brief attacks of sharp, stabbing, or electric-shock pain caused by irritation of nerves at the back of the head. Pain may spread from the upper neck across the scalp and can be associated with tenderness or sensitivity to touch.

\medterm{Occiput} The back part of the skull and the back of the head. It protects the rear of the brain and includes the occipital bone, which joins the skull base and neck.
\medterm{Occiput Posterior Position} A position during labor in which the back of the baby’s head faces the pregnant person’s back. It can cause back pain or slow descent, although many babies rotate to a better position during labor.

\medterm{Occiput Transverse Position} A fetal head position in which the back of the baby’s head faces the pregnant person’s side. It often rotates during labor; if it persists, descent may be slow and assisted delivery or cesarean birth may be needed.

\medterm{Occlude} To block or close a passage, vessel, airway, duct, or opening. An occlusion may be caused by a clot, stone, swelling, scar, tumor, foreign object, or deliberate medical treatment.

\medterm{Occlusal Adjustment} Selective reshaping of tooth surfaces to improve how the upper and lower teeth meet and reduce abnormal contact or bite interference.

\medterm{Occlusal Guard} A removable device worn over the teeth to protect them from grinding, clenching, or injury. A night guard may reduce tooth wear during sleep, while a sports mouthguard protects against blows; the device should fit properly and does not treat every cause of jaw pain.

\medterm{Occlusal Splint} A removable dental appliance that changes or protects tooth contact and may reduce effects of grinding, jaw-joint strain, or selected bite problems.

\medterm{Occlusion} Complete blockage of a blood vessel or body passage. It may be caused by a blood clot, plaque, pressure, swelling, or an intentional medical procedure; the effects depend on the location and can include loss of blood flow, pain, or organ damage.

\medterm{Occlusive Vascular Disease} Disease in which narrowing or blockage of a blood vessel reduces or stops blood flow to tissues.

\medterm{Occult} Hidden or not readily visible or detectable on routine examination. An occult disease or bleeding may be found only with imaging, microscopy, or laboratory testing.

\medterm{Occult Blood} Blood that is not visible to the naked eye and is detected by laboratory testing. In stool, it may be detected by a chemical guaiac test or an immunochemical test; the method and specimen determine the collection instructions and possible dietary or medicine effects. A positive stool test indicates blood is present but does not establish its cause.

\medterm{Occult Blood Screen} Alternate index label; see \textit{Fecal occult blood test}.

\medterm{Occult Gastrointestinal Bleeding} Slow or intermittent bleeding somewhere in the digestive tract that cannot be seen in the stool. It may cause iron-deficiency anemia, tiredness, or a positive stool blood test and requires evaluation to find the source.

\medterm{Occult Primary Cancer} Alternate index label; see \textit{Carcinoma of unknown primary}.

\medterm{Occupational Asthma} Asthma caused or worsened by substances at work, such as flour dust, latex, chemicals, or fumes. Symptoms may improve on days away from work and worsen during exposure; careful exposure history and breathing tests help confirm it.

\medterm{Occupational COPD} Chronic obstructive pulmonary disease caused or worsened by repeated workplace exposure to dust, fumes, chemicals, smoke, or other inhaled hazards. It causes lasting airflow limitation, cough, phlegm, and breathlessness; smoking and other risks may also contribute.

\medterm{Occupational Dermatoses} Skin diseases caused or worsened by workplace conditions, such as repeated wet work, chemicals, gloves, dust, plants, or friction. Common problems include irritant or allergic contact dermatitis, often affecting the hands; reducing exposure and protecting the skin help prevent recurrence.

\medterm{Occupational Disease} An illness caused or significantly worsened by workplace exposures or conditions. Examples include asthma from dust, hearing loss from noise, skin disease from chemicals, and lung disease from harmful particles.
\medterm{Occupational Lung Disease} Lung disease caused or worsened by dusts, fumes, gases, vapors, fibers, or other exposures at work. Examples include occupational asthma, pneumoconiosis, hypersensitivity pneumonitis, and asbestos-related disease.
\medterm{Occupational Exposure} Contact with a potentially harmful substance, energy source, or condition during work. Examples include radiation, chemicals, dust, noise, heat, infection, and repetitive strain. Risk depends on the hazard, amount, duration, and protective controls; monitoring and prevention are part of occupational health.

\medterm{Occupational Exposure Limit} The highest recommended or legally permitted workplace exposure to a hazardous substance over a specified time. The limit varies by substance, country, route of exposure, and whether it applies to a brief peak or an average work shift.

\medterm{Occupational Health} The medical field focused on preventing work-related illness, injury, and harmful exposure while supporting workers' health.
\medterm{Occupational Hearing Loss} Hearing damage caused by repeated or prolonged workplace noise or substances that harm the ear. It is usually permanent and may cause difficulty understanding speech, especially in background noise; reducing exposure and using effective hearing protection help prevent it.

\medterm{Occupational Rehabilitation} Assessment, treatment, and workplace support that help a person return to work or remain employed after illness, injury, or disability. It may include graded activity, skills training, equipment, job changes, and coordination with employers and clinicians.

\medterm{Occupational Safety And Health} Rules and practices that protect workers from injuries, illness, dangerous exposures, and unsafe conditions. Measures include training, protective equipment, ventilation, safe procedures, health monitoring, and reporting or controlling workplace hazards.
\medterm{Occupational Therapist} A healthcare professional who helps people perform daily activities through rehabilitation, adaptive techniques, environmental changes, and assistive equipment.

\medterm{Occupational Therapy} Rehabilitation that helps a person perform daily activities, work, school, and other meaningful tasks after illness, injury, or disability. It may use exercises, practice, adaptive equipment, changes to the environment, and strategies for movement, thinking, or sensory problems.

\medterm{Occupational Therapy Assessment} An evaluation of how illness, injury, disability, or development affects daily activities, school, work, and self-care. The therapist identifies strengths and barriers and recommends exercises, skills practice, equipment, or environmental changes.

\medterm{OCD} Obsessive-compulsive disorder, a condition involving recurrent unwanted thoughts or urges and/or repetitive behaviors or mental acts that cause distress or interfere with daily life.

\synonyms
Obsessive Compulsive Disorder (OCD)

\medterm{Ochre} A yellow-brown color or earth pigment. In medicine, ochre discoloration can describe yellow-brown skin or tissue changes caused by pigment, blood breakdown, chronic congestion, or other conditions.


\medterm{Ochronosis} Blue-black or brown discoloration of cartilage, tendons, skin, or the whites of the eyes caused by buildup of a chemical called homogentisic acid. It is most often linked to the inherited disorder alkaptonuria and may cause arthritis.

\medterm{OCP} OCP commonly means oral contraceptive pill, a medicine containing estrogen, progestin, or both to prevent pregnancy and sometimes treat menstrual disorders.


\medterm{Octreoscan} Nuclear-medicine imaging using indium-111 octreotide, which binds to somatostatin receptors. It was used to locate some neuroendocrine tumors and their spread; newer receptor-targeted PET scans are often preferred.

\medterm{Octreotide} Octreotide is a somatostatin analogue that suppresses growth hormone, several gastrointestinal hormones, and intestinal secretion; it is used for acromegaly, certain neuroendocrine tumors, and selected secretory diarrhoeas.

\medterm{Ocular} Relating to the eye or vision. The term may describe an eye structure, examination, medicine, injury, infection, or disease affecting sight. Ocular emergencies such as sudden severe pain or sudden loss of vision need immediate eye care.
\medterm{Ocular Albinism} Ocular albinism is an inherited disorder with reduced pigment primarily in the eyes, causing decreased visual acuity, nystagmus, photophobia, and abnormal retinal development.

\medterm{Ocular Fixation} The ability to hold the eyes steadily on a target; abnormal fixation can impair visual acuity, tracking, or alignment.

\medterm{Ocular Histoplasmosis Syndrome} An eye condition linked to an earlier histoplasmosis infection. It can cause scars in the retina and, in some people, abnormal blood vessels that damage central vision.

\medterm{Ocular Hypertension} Higher-than-normal pressure inside the eye without current signs of glaucoma damage. It increases the future risk of glaucoma and needs regular eye checks.

\medterm{Ocular Ischemic Syndrome} Eye findings caused by severely reduced blood flow through the arteries supplying the eye, usually from advanced carotid artery disease. It may cause gradual or sudden vision loss, eye pain, retinal changes, or abnormal new blood vessels.

\medterm{Ocular Melanoma} A cancer that begins in pigment-producing cells of the eye.

\medterm{Ocular Migraine} A migraine-related visual disturbance, often involving flashing lights, zigzag lines, or a temporary blind spot. It may occur with or without headache.

\medterm{Ocular Motility} The ability of the eyes to move together in a coordinated way. Abnormal eye movements or poor alignment can cause double vision, eyestrain, dizziness, or reduced vision and may result from muscle, nerve, brain, or eye disease.

\medterm{Ocular Myasthenia Gravis} A form of myasthenia gravis in which weakness affects only the muscles that move the eyes and lift the eyelids. It causes drooping eyelids or double vision that worsens with use and improves with rest. Symptoms may later spread to swallowing, speech, breathing, or limb muscles, so follow-up is important.

\synonyms
OMG; Ocular MG

\medterm{Ocular Myiasis} Infestation of the eye or nearby tissues by fly larvae. It can cause pain, irritation, tearing, and vision problems and requires prompt eye care.

\medterm{Ocular Orbit} The bony cavity around the eye. It contains the eyeball, eye-movement muscles, nerves, blood vessels, glands, and protective tissues.

\medterm{Ocular Pain} Pain in or around the eye caused by surface irritation, infection, inflammation, trauma, glaucoma, or orbital
disease. Severe pain with vision loss, light sensitivity, nausea, or injury requires urgent
ophthalmic assessment.

\medterm{Ocular Rosacea} Rosacea affecting the eyelids or surface of the eyes, causing burning, redness, a gritty feeling, dryness, or light sensitivity. Persistent inflammation can affect the cornea and may impair the eye surface.
\medterm{Ocular Surface Disease} A group of conditions affecting the eye's surface, eyelids, or tears, including dry eye and eyelid inflammation. Symptoms may include burning, redness, irritation, and blurred or changing vision.

\medterm{Ocular Toxoplasmosis} A retinal infection or inflammation caused by Toxoplasma gondii, which can produce blurred vision, floaters, pain, and recurrent retinal scars.

\medterm{Oculocardiac Reflex} A slowing of the heartbeat or similar body response caused by pressure or pulling on the eye or tissues around it.

\medterm{Oculocerebrorenal Syndrome} A rare inherited disorder, also called Zellweger spectrum disorder, that affects the brain, eyes, kidneys, liver, and other organs. Severe forms begin in newborns and can cause weak muscles, seizures, feeding problems, and major developmental problems.

\synonyms
Zellweger spectrum disorder

\medterm{Oculocutaneous Albinism} Alternate index label; see \textit{Albinism}.

\medterm{Oculogyric Crisis} A prolonged, involuntary upward turning of the eyes, often caused by a medicine or neurologic disorder.
\medterm{Oculomotor Nerve} The third cranial nerve, controlling most eye movements, eyelid lifting, and pupil narrowing. Damage can cause double vision, a drooping eyelid, an outward-turning eye, or an abnormally enlarged pupil.
\medterm{Oculomotor Nerve Palsy} Weakness or paralysis of muscles supplied by the third cranial nerve, causing eyelid drooping, double vision, abnormal eye position, and sometimes a dilated pupil.

\medterm{Oculomotor Nucleus} A group of nerve cells in the midbrain that controls most eye movements, lifts the upper eyelid, and narrows the pupil through the third cranial nerve.

\medterm{Oculopharyngeal Muscular Dystrophy} An inherited muscle disease that usually begins in middle age and gradually weakens the eyelids and throat muscles. Drooping eyelids, difficulty swallowing, and later weakness of the hips and shoulders are typical. A gene test can confirm the diagnosis; swallowing assessment, exercises, and surgery may help prevent complications.

\synonyms
OPMD

\medterm{Oculoplastic Procedures} Operations on the eyelids, tear ducts, eye socket, or nearby facial tissues to restore function, repair injury, remove disease, or change appearance. Risks include bleeding, infection, scarring, dry eyes, and rarely injury affecting vision.

\medterm{Oculosympathetic Palsy} Alternate index label; see \textit{Horner syndrome}.

\medterm{ODD} Alternate index label; see \textit{Oppositional defiant disorder}.

\medterm{Oddi, Sphincter Of} A ring-shaped muscle at the opening where bile and pancreatic juice enter the small intestine. It normally opens during digestion and closes afterward; abnormal tightening or blockage can cause pain or pancreatitis.

\medterm{Odontalgia} Pain coming from a tooth or the tissues around it. Common causes include decay, a cracked tooth, infection, gum disease, or injury; swelling, fever, difficulty swallowing, or severe pain needs prompt dental or medical care.
\medterm{Odontoblast} A cell along the inner part of a tooth that makes dentin, the hard layer beneath enamel. Odontoblasts produce dentin while teeth develop and can make limited repairs after injury.
\medterm{Odontoclast} A cell that breaks down the hard tissues of a tooth. This process helps baby teeth be shed naturally, but unwanted tooth resorption can weaken a tooth and may follow injury, inflammation, pressure, or other disease.
\medterm{Odontodysplasia} Odontodysplasia is abnormal development of tooth enamel and dentin, producing thin, weak, discolored teeth that may erupt late or be easily damaged.

\medterm{Odontogenesis} The process by which teeth form and develop before and after birth. Problems during development can cause missing teeth, weak enamel, abnormal dentin, or changes in tooth shape and eruption.
\medterm{Odontogenic Cutaneous Fistula} An abnormal tunnel from an infected tooth to the skin, usually on the chin, cheek, or jaw. It may appear as a small draining opening; treating the tooth infection usually closes the tunnel.

\medterm{Odontogenic Infection} An infection arising from a tooth or the tissues around it, usually caused by dental caries, pulp infection, gum disease, or trauma. It may form an abscess and spread into the jaw, face, neck, or bloodstream.

\medterm{Odontogenic Cyst} A fluid-filled sac arising from tissues involved in tooth development, usually within the jaw. It may be painless or cause swelling, tooth displacement, infection, or bone damage.
\medterm{Odontogenic Fibroma} An odontogenic fibroma is a rare benign tumor arising from tooth-forming connective tissue, usually presenting as a jaw lesion.

\medterm{Odontogenic Keratocyst} A cyst that develops in the jaw from tooth-forming tissue. It may grow into bone without causing symptoms, can recur after removal, and is sometimes associated with an inherited skin-cancer syndrome; imaging and tissue examination guide treatment.

\medterm{Odontogenic Neoplasm} An odontogenic neoplasm is a growth arising from tissues involved in tooth formation, usually in the jaw; it may be benign or malignant.

\medterm{Odontogenic Tumors And Cysts} Alternate index label; see \textit{Jaw tumors and cysts}.

\medterm{Odontoid Process} A peg-like projection from the second neck vertebra that fits into the first neck vertebra and allows the head to turn from side to side.

\medterm{Odontoma} A noncancerous overgrowth of tooth-forming tissue, usually found in the jaw and often discovered when a permanent tooth fails to erupt. It may look like several small teeth or a solid calcified mass and is usually removed if it blocks eruption.

\medterm{Odontopathy} Any disease or abnormality affecting a tooth or the tissues supporting it. Examples include tooth decay, gum disease, infection, abnormal tooth development, wear, and injury. Untreated dental infection can spread to surrounding tissues and occasionally to the bloodstream.


\medterm{Odynophagia} Pain when swallowing, often felt as burning, squeezing, or a sticking sensation in the throat or chest. Causes include infection, reflux, ulcers, medicines, injury, and other swallowing-tube disorders.

\medterm{Oedema} British spelling of edema: abnormal fluid accumulation in tissues that causes swelling. Causes include injury, infection, venous disease, heart or kidney failure, medicines, and lymphatic blockage.

\medterm{Oesophageal} Relating to the oesophagus, the muscular tube carrying swallowed material from the pharynx to the stomach.
\medterm{Oesophageal Atresia} A birth defect in which the swallowing tube ends blindly instead of connecting normally to the stomach. A baby may drool, cough, choke, or turn blue during feeds; urgent specialist care and surgery are required, especially when an abnormal connection to the airway is also present.
\medterm{Oesophageal Ulcer} An open sore in the swallowing tube between the throat and stomach. Reflux, infection, medicines, corrosive injury, inflammation, or cancer may cause it, leading to painful swallowing, chest pain, bleeding, or difficulty swallowing.

\medterm{Oesophageal Varices} Enlarged fragile veins in the wall of the swallowing tube, usually caused by high pressure in the liver’s circulation. They may cause no symptoms until they rupture, leading to vomiting blood or black stools and a life-threatening emergency.
\medterm{Oesophagitis} Inflammation of the swallowing tube between the throat and stomach. Reflux, infections, medicines, allergies, radiation, or corrosive substances can cause it, leading to heartburn, painful or difficult swallowing, chest pain, or bleeding.

\medterm{Oesophagoscopy} Examination of the food pipe using a thin viewing tube passed through the mouth. It can investigate swallowing difficulty, pain, bleeding, or abnormal tissue, and can remove a foreign body, take samples, stop bleeding, or widen a narrowing. Sedation or anesthesia is usually used.
\medterm{Oesophagus} The muscular tube that transports swallowed food and liquid from the pharynx to the stomach.
\medterm{Oesophagus, Diseases Of} Disorders affecting the swallowing tube between the throat and stomach. They include reflux, inflammation, narrowing, movement problems, cancer, infection, and birth defects, often causing swallowing difficulty, pain, heartburn, or regurgitation.
\medterm{Oestradiol} British spelling of estradiol, the main naturally occurring estrogen during the reproductive years. It supports reproductive development, the menstrual cycle, bone health, and other body functions.
\medterm{Oestriol} A form of estrogen made in large amounts during pregnancy by the placenta using substances from the fetus and pregnant person. Its level can be measured in some prenatal screening tests, but it is not interpreted alone to diagnose a problem.
\medterm{Oestrogen} A group of hormones important for reproductive development, menstrual cycles, pregnancy, bone strength, and other body functions. They are made mainly by the ovaries before menopause and by other tissues, including the placenta during pregnancy.

\medterm{Oestrogen Replacement Therapy} Treatment with estrogen, sometimes combined with a progestogen, to relieve symptoms caused by low estrogen such as hot flushes and vaginal dryness. Benefits and risks depend on age, medical history, uterus status, dose, and route, so treatment is individualized.

\medterm{Oestrogens} A group of steroid hormones involved in reproductive development, menstrual cycling, pregnancy, bone, and other physiological functions.
\medterm{Oestrone} A form of estrogen made by the ovaries and by conversion of other hormones in body tissues. It becomes the main estrogen after menopause and affects reproductive tissues, bones, blood vessels, and metabolism.
\medterm{Ofatumumab} A monoclonal antibody against CD20 used for selected blood cancers and relapsing multiple sclerosis; it can increase infection risk and reduce B cells.

\medterm{Off-Label} Use of an approved medicine for an indication, dose, route, or population not specified in its regulatory approval.
\medterm{Off-Label Use} Use of an approved medicine for a condition, dose, route, or age group not specifically included in its official approval. A clinician may recommend it when evidence and expected benefits support its use; the decision should be discussed with the patient.

\medterm{Offspring} The young or descendants produced by one organism or by a pair. In medicine and genetics, offspring may be studied to understand inherited traits, pregnancy outcomes, developmental conditions, and disease risk.

\medterm{Ofloxacin} A fluoroquinolone antibiotic used for selected susceptible bacterial infections. It can cause serious tendon, nerve, heart-rhythm, and other adverse effects, so it should be used only when clinically appropriate.
\medterm{Ogilvie Syndrome} A dangerous widening of the large intestine without a physical blockage, usually in a seriously ill or recently operated patient. Abdominal swelling, pain, nausea, or failure to pass stool may occur, and the bowel can perforate if untreated.

\medterm{Ogilvie’S Syndrome} Another spelling of Ogilvie syndrome: severe widening of the large intestine without a physical blockage, often after surgery or during serious illness. It needs medical monitoring because pressure can reduce blood flow or cause a bowel perforation.

\medterm{Ohm} The standard unit of electrical resistance, equal to one volt per ampere. In medical equipment, resistance affects how electrical current flows through devices, sensors, wires, and body tissues.
\medterm{OHSS} Alternate index label; see \textit{Ovarian hyperstimulation syndrome}.


\medterm{Oidium} A fragment or spore-like fungal element formed by fragmentation of a hypha; the term is also used in older fungal classifications.
\medterm{Oil} A hydrophobic liquid used as a nutrient, solvent, emollient, or medicinal preparation; safety depends on the specific oil.
\medterm{Oil-Based Paint Poisoning} Harm from swallowing or inhaling oil-based paint or its solvents. It may irritate the lungs, cause dizziness or drowsiness, and cause chemical pneumonia if liquid enters the lungs; do not induce vomiting.

\medterm{Oily Hair} Hair that appears greasy because the scalp produces excess oil or because of product buildup. It may occur with infrequent washing, hormonal changes, or some skin conditions and is not poisoning.

\medterm{Oily Skin} Skin that produces more oil than usual, often causing shine, enlarged-looking pores, or acne. Gentle cleansing and non-comedogenic products may help; persistent acne or irritation may need medical assessment.

\medterm{Ointment} A thick, semisolid preparation applied to the skin, inside the nose, or other body surfaces. Its oily base helps protect the area and keep it moist; some ointments also contain medicines that treat infection, inflammation, pain, or other conditions.

\medterm{Ointment Bases} Semisolid materials that carry medicines applied to the skin or moist body surfaces. They affect how easily a product spreads, how much moisture it retains, how much medicine enters tissue, and how long it remains stable.

\medterm{Okadaic Acid} A toxin made by certain marine organisms that can accumulate in shellfish. Eating contaminated shellfish may cause diarrhoea and vomiting; in laboratories it is used to study enzymes that remove phosphate groups from proteins.
\medterm{Okihiro Syndrome} A rare inherited condition caused by changes in the SALL4 gene. It may combine abnormalities of the thumb or forearm with difficulty moving the eyes, and may also affect hearing, the heart, kidneys, or spine.

\medterm{Olanzapine} Olanzapine is an atypical antipsychotic that blocks dopamine and serotonin receptors and treats schizophrenia and bipolar disorder; weight gain, diabetes, dyslipidaemia, sedation, and movement disorders are important risks.

\medterm{Olaparib} Olaparib is a PARP inhibitor used for selected cancers with BRCA or other homologous-recombination repair defects; anaemia, nausea, fatigue, myelosuppression, and rare myelodysplastic syndrome or acute myeloid leukaemia may occur.


\medterm{Old Age} Later life; health assessment should consider function, comorbidity, cognition, and social circumstances rather than age alone.
\medterm{Oleander Poisoning} Poisoning from eating oleander, which contains heart-active toxins. Nausea, vomiting, dizziness, confusion, slow or irregular heartbeat, fainting, or seizures require immediate emergency care.


\medterm{Olecranon} The prominent bony process at the proximal end of the ulna forming the point of the elbow.
\medterm{Oleic Acid} A monounsaturated fat found in olive oil, nuts, seeds, avocados, and many other foods. Replacing some saturated fats with it may support healthier blood cholesterol levels as part of a balanced diet.
\medterm{Oleum} A very corrosive mixture of sulfur trioxide and sulfuric acid, used industrially and dangerous to skin, eyes, and lungs.
\medterm{Olfaction} The sense of smell. Odor molecules stimulate special cells in the upper part of the nose, which send signals to the brain to identify and perceive odors.

\medterm{Olfactory} Relating to the sense of smell or the nerves and brain pathways that process odors. Problems in this system can cause reduced smell, loss of smell, distorted smells, or unusual odor sensations.
\medterm{Olfactory Bulb} The olfactory bulb is a forebrain structure that receives signals from nasal smell receptors and begins processing odor information.

\medterm{Olfactory Cortex} The olfactory cortex is a group of brain regions that interprets smell signals and links them with memory, emotion, and behavior.

\medterm{Olfactory Gyrus} The olfactory gyrus is a cortical region associated with processing olfactory information and includes areas near the medial temporal and frontal lobes.

\medterm{Olfactory Learning} Olfactory learning is acquisition or modification of behavior or memory through exposure to odors and their associations.

\medterm{Olfactory Mucosa} The olfactory mucosa is specialized nasal lining containing receptor neurons, supporting cells, and glands that detect odor molecules.

\medterm{Olfactory Nerve} The first cranial nerve, which carries smell signals from specialized cells high in the nose to the brain. Injury or disease can reduce or eliminate the sense of smell or distort odors.

\medterm{Olfactory Neuroblastoma} A rare malignant tumor arising in the upper nasal cavity near the olfactory nerves, often causing nasal blockage, nosebleeds, or reduced smell.

\medterm{Olfactory Tract} The olfactory tract is a bundle of nerve fibers carrying smell information from the olfactory bulb to cerebral regions involved in odor perception.

\medterm{Olfactory Tubercle} The olfactory tubercle is a ventral forebrain region involved in odor processing, reward, motivation, and integration with limbic circuits.



\medterm{Oligoasthenoteratozoospermia} A semen abnormality in which sperm numbers are low, movement is reduced, and fewer sperm have a typical shape. It can reduce fertility and may result from testicular disease, hormone problems, heat, medicines, smoking, or an unknown cause. Repeat semen testing and evaluation help guide treatment or assisted reproduction.

\synonyms
OAT Syndrome; OAT

\medterm{Oligoclonal Bands} Groups of immune proteins found in blood or the fluid surrounding the brain and spinal cord. Bands found only in that fluid suggest immune activity inside the nervous system, but they can occur in multiple sclerosis, infections, and other inflammatory conditions.

\medterm{Oligocythaemia} An unusually low number of blood cells, especially red blood cells. The uncommon term may describe anemia or another low blood-cell count, so the specific cell type and cause must be identified with a blood test.
\medterm{Oligodendrocyte} A supporting cell in the central nervous system that forms myelin around axons. Myelin helps nerve signals travel quickly; oligodendrocytes can be damaged in disorders such as multiple sclerosis.
\medterm{Oligodendroglia} Oligodendroglia are glial cells of the central nervous system that produce myelin around axons and support nerve conduction.

\medterm{Oligodendroglial Neoplasm} A brain tumor made of cells that resemble oligodendrocytes, the cells that form the insulating coating around nerve fibers. Diagnosis uses tissue examination and genetic testing, and the tumor’s behavior depends on its grade and molecular features.

\medterm{Oligodendroglioma} A brain tumor defined by characteristic changes in the IDH gene and loss of parts of chromosomes 1 and 19. It usually begins in the cerebral hemispheres and may cause seizures, headache, or other neurologic symptoms; grade affects its behavior and treatment.

\medterm{Oligodontia} A condition present from birth in which many permanent teeth fail to develop, usually more than six excluding wisdom teeth. It may affect chewing, speech, appearance, and jaw development, requiring dental replacement or orthodontic care.
\medterm{Oligogenic Inheritance} A pattern in which changes in a small number of genes work together to influence a disease or characteristic. The combination and strength of the changes can make symptoms differ among affected people.

\medterm{Oligohydramnios} Too little amniotic fluid around the baby during pregnancy. It may result from leaking membranes, placental problems, dehydration, a pregnancy that continues beyond the due date, or problems affecting the baby's kidneys.

\medterm{Oligomenorrhea} Oligomenorrhea is infrequent or widely spaced menstrual bleeding, often from anovulation, polycystic ovary syndrome, endocrine disease, undernutrition, excessive exercise, stress, or perimenopause.

\medterm{Oligomenorrhoea} Menstrual periods that occur less often than expected, usually with cycles longer than about 35 days. Causes include pregnancy, polycystic ovary syndrome, thyroid disease, stress, low body weight, intense exercise, and some medicines; evaluation depends on age and symptoms.
\medterm{Oligometastasis} Cancer that has spread to only a small number of separate sites, often no more than five. In selected people, surgery, focused radiation, or another local treatment may control these deposits for a long time.


\medterm{Oligonucleotide} A short piece of DNA or RNA made in the laboratory or found in cells. It can be used to detect genetic material, block or change gene messages, or help read and study genes.

\medterm{Oligopeptide} A short chain of amino acids, usually containing only a few to about 20 building blocks. Oligopeptides can act as hormones, signaling molecules, antimicrobial substances, or parts of larger proteins.

\medterm{Oligosaccharide} A short chain of sugar molecules joined together. Oligosaccharides form parts of cell-surface molecules, affect how cells recognize one another, and occur naturally in foods where some are used by gut bacteria.

\medterm{Oligospermia} A lower-than-expected sperm concentration in semen. It can reduce fertility and may result from hormone problems, infections, medicines, heat, toxins, varicoceles, genetic conditions, or problems with sperm production.
\medterm{Oliguria} Abnormally low urine production. It may result from dehydration, low blood flow, blockage, kidney disease, or medicines and can be an early sign of sudden kidney injury.

\medterm{Olivary} Relating to the olive-shaped inferior olivary nucleus of the medulla or to an olive-like structure.
\medterm{Olivary Nucleus} The olivary nucleus is a brainstem structure that sends timing and motor-learning signals to the cerebellum through climbing fibers.

\medterm{Olive Oil} A plant oil rich in monounsaturated fatty acids, used as food and in some topical preparations.
\medterm{Olive Pontocerebellar Atrophy} A group of progressive brain disorders that damage areas controlling balance and coordination. They may cause worsening unsteadiness, clumsy movements, slurred speech, stiffness, or other neurologic problems.

\medterm{Olivopontocerebellar Atrophy} Olivopontocerebellar atrophy is a degenerative condition affecting the inferior olives, pons, and cerebellum, causing progressive ataxia and other neurologic dysfunction.

\medterm{Ollier's Disease} A nonhereditary skeletal disorder with multiple enchondromas, causing deformity, limb-length discrepancy, and an increased risk of malignancy.
\medterm{Olmesartan Medoxomil} An oral prodrug converted to olmesartan, an angiotensin II receptor blocker used to treat high blood pressure.

\medterm{Olopatadine Hydrochloride} Olopatadine hydrochloride is an antihistamine used in eye drops or nasal spray to treat allergic conjunctivitis or rhinitis.

\medterm{OLP} Alternate index label; see \textit{Oral lichen planus}.

\medterm{Olsenella} Olsenella is a genus of anaerobic bacteria found mainly in the mouth and digestive tract; human infection is uncommon.

\medterm{Olsenella Uli} An anaerobic gram-positive bacterium found in the mouth and associated with periodontal disease; invasive infection is rare.


\medterm{Omacetaxine Mepesuccinate} Omacetaxine mepesuccinate is a protein-synthesis inhibitor used for selected chronic myeloid leukemia resistant or intolerant to tyrosine-kinase inhibitors.

\medterm{Omalizumab} Omalizumab is a monoclonal antibody that binds circulating IgE and reduces allergic inflammation; it is used for selected allergic asthma, chronic urticaria, nasal polyps, and food allergy, with rare anaphylaxis.

\medterm{Omasum} The omasum is the third compartment of a ruminant stomach, absorbing water and small particles before contents reach the abomasum.

\medterm{Omega-3} A family of polyunsaturated fats found in fish, flaxseed, walnuts, and some supplements, supporting heart and brain health.


\medterm{Omega-3 Fatty Acids} A family of fats found in oily fish, seafood, some seeds, nuts, and plant oils. EPA and DHA support cell and eye function, while ALA is a plant form; supplements are not appropriate for everyone and can interact with medicines.

\medterm{Omega-6} A family of essential polyunsaturated fats, found in vegetable oils, nuts, seeds, and many prepared foods. The body uses them for cell membranes and signaling. Linoleic acid is the main dietary form; replacing saturated fats with unsaturated fats generally supports heart health.

\medterm{Omega-6 Fatty Acids} Essential polyunsaturated fats found in plant oils, nuts, seeds, and many prepared foods. The body needs them for cell membranes and signaling, but overall health depends on balanced intake with other fats.

\synonyms
n-6 fatty acids

\medterm{Omega-N-Methylarginine} Omega-N-methylarginine is a methylated arginine compound that can inhibit nitric-oxide synthase and affect vascular signaling.
\medterm{Omenn Syndrome} A severe inherited immune deficiency that usually begins in infancy. Babies may develop widespread red, peeling skin, swollen lymph nodes, enlarged liver or spleen, diarrhea, poor growth, and repeated serious infections.

\medterm{Omentectomy} Removal of part or all of the greater omentum, commonly performed during staging and
cytoreduction for ovarian and other intra-abdominal cancers.

\medterm{Omentum} A fold of fatty tissue inside the abdomen that hangs from the stomach and partly covers the intestines. It contains blood vessels and immune cells and can help limit the spread of infection or inflammation.

\medterm{Omeprazole} Omeprazole is a proton-pump inhibitor that irreversibly suppresses gastric H+/K+-ATPase; it treats reflux and peptic-ulcer disease and supports Helicobacter pylori eradication, while prolonged use may affect magnesium, B12, bone, and infection risk.

\medterm{Ommaya Reservoir} A surgically implanted reservoir connected to a catheter in a brain ventricle, used to sample cerebrospinal fluid or deliver selected medicines.

\medterm{Omphalitis} Infection of a newborn’s umbilical stump, usually during the first weeks of life. Redness spreading onto the skin, swelling, pus, bad smell, fever, poor feeding, or unusual sleepiness requires urgent treatment because infection can spread rapidly.

\medterm{Omphalocele} A birth defect in which abdominal organs protrude through an opening at the navel into a protective sac. Other chromosome or organ abnormalities may occur, so newborn assessment and imaging are needed.

\medterm{Omphalocele Repair} Treatment to place an infant's abdominal organs back inside the abdomen and close the opening in the abdominal wall. Repair may be immediate or staged according to the size of the defect and the baby's breathing and circulation.

\medterm{Onchocerca Volvulus} A parasitic roundworm spread by repeated bites of infected blackflies that breed near fast-flowing rivers. Adult worms live in lumps under the skin and release young worms that can migrate through the skin and eyes.

\medterm{Onchocerciasis} A parasitic infection caused by Onchocerca volvulus and spread by infected blackflies. It can cause intensely itchy skin, nodules, skin changes, and eye inflammation that may lead to permanent visual impairment or blindness.

\medterm{Oncocytic Neoplasm} A growth made of cells containing unusually many energy-producing structures, giving them a granular appearance under a microscope. It may be noncancerous or cancerous, depending on the organ and the tumor’s other features.

\medterm{Oncocytoma} A usually noncancerous tumor made of cells containing many energy-producing structures. It most often develops in the kidney but can occur in other organs; testing is needed to distinguish it from cancer.

\medterm{Oncogene} A gene that can make cells grow or survive without normal control when it is changed, overactive, or present in too many copies. Oncogene changes can contribute to the development and progression of cancer.

\medterm{Oncogenic} Capable of causing or promoting cancer. An oncogenic change, virus, gene, or chemical can increase abnormal cell growth or reduce normal control of cell division.

\medterm{Oncogenic Virus} A virus capable of contributing to cancer by altering cell growth, suppressing tumor defenses, or integrating or maintaining oncogenic genetic material.

\medterm{Oncogenicity} The ability of a substance, infection, genetic change, or other factor to contribute to the development of cancer.

\medterm{Oncologic Emergency} A cancer-related problem that can quickly threaten life, organ function, or the nervous system and needs urgent treatment. Examples include spinal-cord pressure, blocked major chest veins, dangerously high calcium, tumor breakdown, airway blockage, and infection during low white-cell counts.

\medterm{Oncologist} A doctor who specializes in cancer care. Medical oncologists use medicines, radiation oncologists use radiation, and surgical oncologists operate; many people receive care from more than one type of oncologist.

\medterm{Oncology} The medical field concerned with preventing, finding, and treating cancer and supporting people after treatment. Oncologists may use surgery, radiation, medicines, or combinations of treatments and work with other professionals to provide symptom and emotional support.
\medterm{Oncology Nursing} Oncology nursing provides cancer-related assessment, treatment support, symptom control, education, coordination, survivorship, and end-of-life care.
\medterm{Oncology PET} A positron-emission tomography scan, often combined with CT, that uses a small radioactive tracer to show active cancer in the body. It may help assess spread, treatment response, or recurrence, but inflammation can also appear active.

\medterm{Oncology Rehabilitation} Healthcare that helps people prevent or manage cancer-related problems before, during, or after treatment. It may address fatigue, pain, weakness, nerve symptoms, swelling, movement, nutrition, work, daily activities, and emotional wellbeing.

\medterm{Oncolytic Virotherapy} Oncolytic virotherapy uses modified viruses that preferentially infect and destroy cancer cells and may stimulate an immune response against the tumor.

\medterm{Oncolytic Virus} A virus selected or modified to infect and destroy cancer cells more readily than normal cells. It may also stimulate the immune system to attack the tumor; use is limited to selected cancers and requires specialist treatment.

\medterm{Oncolytic Viruses} Viruses selected or engineered to infect and destroy cancer cells while stimulating an immune response against the tumor.

\medterm{Oncoprotein} A protein that promotes abnormal cell growth or survival and can contribute to cancer. It may result from a cancer-causing gene, a mutation, a chromosome change, or infection with certain viruses.

\medterm{Oncosis} A form of cell injury in which cells swell and may die, usually because severe lack of oxygen or another injury disrupts energy production and the movement of salt and water across the cell membrane.


\medterm{Oncotic Pressure} The pulling force created by proteins in the blood, mainly albumin, that helps keep fluid inside blood vessels. When blood protein is low, fluid may leak into tissues and cause swelling.

\medterm{Oncotoxic} Harmful to cancer cells or relating to substances that damage tumors. An oncotoxic treatment may also affect healthy cells, so its benefits, side effects, dose, and safety require careful medical monitoring.
\medterm{Ondansetron} Ondansetron is a selective 5-HT3-receptor antagonist used to prevent nausea and vomiting after chemotherapy, radiotherapy, surgery, or other triggers; constipation, headache, and QT prolongation can occur.

\medterm{Ondansetron Hydrochloride} The hydrochloride salt of ondansetron, a serotonin 5-HT3 receptor blocker used to prevent nausea and vomiting.

\medterm{Ondine's Curse} A rare disorder in which automatic breathing is impaired, especially during sleep; also called congenital central hypoventilation syndrome.

\synonyms
congenital central hypoventilation syndrome.

\medterm{Ondine’S Curse} A rare disorder in which automatic breathing becomes inadequate, especially during sleep, because brain signals fail to drive breathing normally. It can cause dangerously low oxygen and may require assisted ventilation.
\synonyms
Congenital central hypoventilation syndrome

\medterm{One-Hour Glucose Challenge Test} A screening test for diabetes during pregnancy, usually performed at 24–28 weeks. After drinking a measured glucose solution, blood sugar is checked one hour later. A high result does not diagnose diabetes but usually leads to a longer confirmatory test.

\medterm{Oneirophrenia} A dreamlike state with confusion or impaired reality testing; the term is uncommon and does not represent a specific modern psychiatric diagnosis.

\medterm{Onset} Onset is the time or manner in which a symptom, disease, or physiologic change begins, such as sudden, gradual, acute, or insidious onset.

\medterm{Onychia} Inflammation or infection of the nail or nail-producing tissue. It can cause redness, swelling, pain, pus, or changes in nail growth and may result from bacteria, fungi, injury, or skin disease. Treatment depends on the cause.
\medterm{Onychocryptosis} Onychocryptosis is an ingrown nail, usually at the great toe, in which the nail edge penetrates surrounding skin and causes pain or inflammation.

\medterm{Onychocryptosis (Ingrowing Toenail)} A condition in which the edge of a nail grows into the surrounding skin, causing pain, inflammation, and sometimes infection.

\synonyms
Ingrowing Toenail

\medterm{Onychocryptosis (Ingrown Toenail)} Alternate index label for Ingrown Toenail.

\medterm{Onychogryphosis} Severe thickening and overgrowth of a nail, which becomes curved and may resemble a claw. It is more common in older adults and can follow injury, fungal infection, poor circulation, or difficulty caring for the feet; it may cause pain in shoes.
\medterm{Onycholysis} Separation of the nail plate from the underlying nail bed, beginning distally or
laterally. Causes include trauma, psoriasis, fungal infection, thyroid disease,
medications, and chemical exposure; the underlying cause should be identified.

\medterm{Onychomadesis} Onychomadesis is proximal separation and shedding of a nail caused by temporary arrest of nail-matrix growth, often after infection, systemic illness, trauma, or medication exposure.

\medterm{Onychomycosis} A fungal infection of a nail, more common in toenails than fingernails. The nail may become thick, yellow or white, brittle, crumbly, or separated from the skin. A sample may be tested before prolonged antifungal treatment because other nail diseases look similar.

\medterm{Onychomycosis (Fungal Nails)} Alternate index label for Fungal Nails.

\medterm{Oocyte} An immature egg cell in an ovary. It develops before birth, remains paused until a menstrual cycle, and may mature and be released for possible fertilization.

\medterm{Oocyte Cryopreservation} Freezing unfertilized eggs for possible use in future pregnancy. Hormone medicines stimulate the ovaries, and the eggs are collected through the vagina with a needle before being rapidly frozen. Eggs may be stored before cancer treatment or for other fertility-preservation reasons, but freezing does not guarantee a future pregnancy.

\synonyms
Egg Freezing; Oocyte Vitrification

\medterm{Oocyte Retrieval} A procedure used in assisted reproduction to collect mature eggs from the ovaries. After hormone stimulation, a clinician guides a thin needle through the vagina using ultrasound and gently draws fluid containing the eggs into a container for laboratory fertilization.

\synonyms
Ovum Pick-Up (OPU); Egg Retrieval; Follicular Aspiration

\medterm{Oogenesis} The formation and maturation of egg cells, or ova, in the ovaries. It begins before birth and continues through the reproductive years, with changes during each menstrual cycle.
\medterm{Oogonium} An oogonium is an early female germ cell that proliferates during embryonic development and gives rise to primary oocytes.

\medterm{Oophorectomy} Surgery to remove one or both ovaries. It may treat cancer, a cyst, infection, torsion, or reduce inherited cancer risk; removing both ovaries before natural menopause causes an immediate fall in ovarian hormones.

\medterm{Oophoritis} Inflammation of an ovary, usually associated with pelvic infection or, less commonly, autoimmune disease.
It may cause lower abdominal or pelvic pain, fever, abnormal discharge, or menstrual changes;
severe pain requires assessment to exclude ovarian torsion or abscess.

\medterm{Ooplasm} Ooplasm is the cytoplasm of an egg cell, containing organelles, nutrients, and molecular factors that support fertilization and early embryo development.

\medterm{Opacity} Clouding of a normally clear part of the eye, such as the cornea, lens, or fluid inside the eye. It can blur vision and may result from injury, infection, inflammation, swelling, or age-related change.

\medterm{Open Bite} An open bite is a malocclusion in which upper and lower teeth do not meet in part of the mouth when the jaws are closed.

\medterm{Open Fractures} Broken bones associated with a wound that connects the fracture to the outside air. They have a high risk of infection and tissue damage and require urgent cleaning, antibiotics, and specialist treatment.

\medterm{Open Gallbladder Removal} Surgery to remove the gallbladder through one larger abdominal incision. It may be needed when keyhole surgery is unsafe or cannot be completed; bleeding, infection, bile leakage, and injury to nearby ducts or organs are possible.

\medterm{Open Lung Biopsy} An open lung biopsy removes a small piece of lung through a surgical chest incision for microscopic diagnosis when less invasive tests are insufficient.

\medterm{Open Pleural Biopsy} An open pleural biopsy removes a sample of the lining around the lung through surgery to investigate cancer, infection, or unexplained pleural disease.

\medterm{Open Prostatectomy} Surgery through an abdominal incision to remove the part of an enlarged prostate that blocks urine flow. It can improve severe urinary obstruction but may cause bleeding, infection, urine leakage, or changes in sexual and urinary function.

\medterm{Open Reduction And Internal Fixation} Surgery to expose a broken bone, put the pieces back into position, and hold them with plates, screws, wires, or rods. It is used when a cast or other nonsurgical treatment cannot provide safe alignment or stability.


\medterm{Open Surgery} Surgery performed through a relatively large incision that allows the surgeon to see and reach the affected area directly. It is used when a larger working space is needed, although recovery may be longer than after some keyhole procedures.

\medterm{Open-Angle Glaucoma} A type of glaucoma in which the eye’s drainage angle remains open, but fluid drains too slowly. Pressure may rise and damage the optic nerve, causing gradual, permanent vision loss, often without early symptoms.

\medterm{Open-Book Fracture} A severe pelvic injury in which the front or back pelvic joints separate, making the pelvis widen like an open book. It usually follows major trauma and can cause life-threatening internal bleeding or organ injury.

\medterm{Open-Heart Surgery} Surgery in which the chest is opened to repair or treat the heart or its major blood vessels. Some operations use a heart–lung machine while the heart is temporarily stopped; risks depend on the operation and the person’s health.

\medterm{Operable} Able to be treated with surgery because the disease can be safely reached or removed.
\medterm{Operant Conditioning} Learning in which the likelihood of a behavior changes because of its consequences, such as reinforcement or punishment.

\medterm{Operating Room} An operating room is a controlled sterile or semi-sterile clinical area equipped for surgical procedures, anesthesia, monitoring, and recovery support.

\medterm{Operating Tables} Adjustable tables designed to support and position patients during surgery or other procedures while allowing access and maintaining safety.

\medterm{Operation} A planned surgical procedure used to diagnose, repair, remove, replace, or improve part of the body. It may be performed through an incision or with instruments inserted through small openings, depending on the condition and goal.
\medterm{Operative Dentistry} The branch of dentistry that repairs teeth damaged by decay, fractures, wear, or developmental defects. Treatment may include fillings, bonding, reshaping, and other procedures that restore appearance, strength, and function.
\medterm{Operative Vaginal Delivery} Vaginal birth assisted with a vacuum cup or forceps when the baby's position is suitable and birth needs to be completed promptly. Reasons may include prolonged pushing or concern about the baby's condition; tears, bleeding, or injury to the baby can occur.

\medterm{Operator} A person who performs a procedure or operates equipment; in healthcare, training and authorization are essential.
\medterm{Operculum} A lid, covering, or flap-like structure; in anatomy it may cover an opening or cortical region.

\medterm{Ophthalmia} Inflammation of the eye, especially the conjunctiva or other surface tissues. It may cause redness, discharge, pain, light sensitivity, or blurred vision and can result from infection, allergy, injury, or inflammation.
\medterm{Ophthalmia Neonatorum} Inflammation or infection of a newborn baby's eye during the first weeks of life. It may be caused by bacteria passed during birth, other infections, or chemical irritation; redness, swelling, or pus requires prompt treatment to protect the cornea and vision.

\medterm{Ophthalmic} Relating to the eye. An ophthalmic medicine, drop, ointment, or device is designed for use on or around the eye and must be used as directed to prevent injury or infection.
\medterm{Ophthalmic Emergency} An eye problem that needs immediate assessment to prevent permanent vision loss or serious injury. Examples include sudden vision loss, retinal detachment, severe eye infection, chemical exposure, a penetrating injury, or rapidly painful high eye pressure.

\medterm{Ophthalmic Solution} A sterile liquid medicine made for application to the eye. It must remain clean and should not be shared or used after contamination because contaminated drops can cause serious eye infections.

\medterm{Ophthalmitis} Inflammation of the eye caused by infection, immune disease, injury, or surgery. It may cause redness, pain, light sensitivity, discharge, or blurred vision and needs prompt assessment when severe.

\medterm{Ophthalmologist} A medical doctor trained to diagnose and treat eye disease and vision problems. Ophthalmologists can prescribe medicines and glasses, perform eye procedures and surgery, and manage conditions affecting the eye, optic nerve, and visual system.

\medterm{Ophthalmology} The medical specialty concerned with the eyes, vision, and the parts of the nervous system used for seeing. It includes eye examinations, treatment of eye disease, vision correction, and procedures or surgery when needed.
\medterm{Ophthalmomyiasis} Infestation of the eye by fly larvae. Larvae may remain on the eye surface or enter deeper tissues, causing intense irritation, pain, inflammation, or vision loss; prompt examination is needed for safe removal and treatment.

\medterm{Ophthalmoplegia} Weakness or paralysis of one or more muscles that move the eyes, causing limited eye movement, double vision, or an abnormal eye position. Causes include nerve disease, stroke, diabetes, thyroid disease, injury, and muscle disorders.

\medterm{Ophthalmoscope} A lighted instrument used to examine the back and internal structures of the eye, including the retina, macula, optic nerve, and retinal blood vessels. Direct ophthalmoscopes provide a magnified view through the pupil, while indirect ophthalmoscopes use a separate light and lens to provide a wider view. Findings may relate to diabetes, high blood pressure, glaucoma, infection, inflammation, or other eye and nervous-system disorders.

\medterm{Ophthalmoscopy} Examination of the inside of the eye with a lighted instrument. It allows a clinician to inspect the retina, optic nerve, and blood vessels for changes caused by diabetes, high blood pressure, glaucoma, inflammation, or other disease.

\medterm{Opiate} A drug derived directly from chemicals in the opium poppy, such as morphine or codeine, or a closely related semisynthetic form. Opiates can relieve pain but may cause sleepiness, constipation, dependence, and dangerously slow breathing.

\medterm{Opiate Alkaloid} An opiate alkaloid is a naturally occurring alkaloid from the opium poppy, such as morphine or codeine, with opioid analgesic and dependence potential.

\medterm{Opiate And Opioid Withdrawal} Opioid withdrawal is a predictable syndrome after reducing or stopping opioid exposure, producing anxiety, sweating, muscle aches, diarrhea, abdominal cramps, nausea, and autonomic symptoms.

\medterm{Opiate Assay} A laboratory test that detects or measures opiates or their breakdown products in urine, blood, saliva, or another specimen. Results may support clinical care or toxicology testing, but some synthetic opioids require separate tests.
\medterm{Opiate Overdose} A life-threatening amount of an opioid medicine or drug that slows brain and breathing activity. Extreme sleepiness, pinpoint pupils, slow or absent breathing, blue lips, or inability to wake require emergency help and naloxone when available.

\medterm{Opioid} A natural or synthetic substance that acts on opioid receptors to relieve pain and may cause sedation or dependence.

\medterm{Opioid Agonist Treatment} Treatment for opioid use disorder using a prescribed medicine, usually methadone or buprenorphine, that activates opioid receptors in a controlled way. It reduces withdrawal and cravings, lowers illicit opioid use and overdose risk, and works best with ongoing medical support.

\medterm{Opioid Intoxication} Opioid intoxication can cause euphoria, drowsiness, slowed breathing, pinpoint pupils, impaired coordination, coma, and death; respiratory depression requires immediate emergency treatment.

\medterm{Opioid Overdose} A medical emergency characterized by the classic triad of respiratory depression, miosis
(pinpoint pupils), and decreased level of consciousness. It can progress to respiratory
arrest, cardiovascular collapse, and death.

\medterm{Opioid Poisoning} Alternate index label; see \textit{Opioid overdose}.

\medterm{Opioid Receptor} A cell-surface protein activated or blocked by opioids and naturally produced endorphins. Mu, kappa, and delta receptors influence pain, breathing, mood, bowel movement, and reward; excessive mu-receptor activation can suppress breathing.

\medterm{Opioid Reversal} Emergency treatment of opioid toxicity with naloxone, a competitive opioid-receptor
antagonist. It is titrated to restore adequate breathing, with repeated doses or monitoring sometimes needed because
its effect may wear off before the opioid.

\medterm{Opioid Toxicity} Toxic effects of opioid agonists, classically respiratory depression, central nervous
system depression, and miosis. Management prioritizes ventilation and airway support;
naloxone is titrated to restore adequate breathing while recognizing the risk of
recurrent toxicity and withdrawal.

\medterm{Opioid Toxidrome} A poisoning pattern caused by opioid excess, characterized by respiratory depression, reduced consciousness, and small pupils; naloxone and ventilatory support may be required.

\medterm{Opioid Use Disorder} A pattern of opioid use that causes loss of control, cravings, continued use despite harm, or problems at home, work, or school. It can occur with prescribed medicines or illicit opioids and is treatable with medication, counseling, and support.

\medterm{Opioid Withdrawal} A characteristic group of symptoms occurring after abrupt reduction or cessation of prolonged opioid use, including anxiety, muscle aches, sweating, abdominal cramps, diarrhea, nausea, and insomnia.

\medterm{Opisthorchis Felineus} A liver fluke acquired by eating raw or undercooked freshwater fish; chronic infection can cause bile-duct inflammation and increase cholangiocarcinoma risk.

\medterm{Opisthotonos} Opisthotonos is severe backward arching of the head, neck, and spine from sustained muscle spasm, classically associated with tetanus, meningitis, strychnine poisoning, or severe neurologic injury.

\medterm{Opium} The dried latex of Papaver somniferum containing opioid alkaloids; it can cause dependence and respiratory depression.
\medterm{Opponens} A muscle that moves a body part toward an opposing surface, especially the thumb across the palm.
\medterm{Opportunistic Infection} An infection that causes illness mainly when the immune system is weakened or a normal protective barrier is damaged. Causes may include cancer treatment, HIV, transplantation, severe illness, or prolonged use of immune-suppressing medicines.

\medterm{Opposition} Movement of the thumb across the palm so it can touch the fingertips. It is important for grasping and depends on several small movements at the joint where the thumb meets the wrist.

\medterm{Oppositional Defiant Disorder} A childhood or adolescent disorder involving a persistent pattern, lasting at least six months, of angry or irritable mood, argumentative or defiant behavior, or vindictiveness toward an authority figure or other person. The pattern causes clinically significant distress or impairment and is not limited to conflicts with siblings.


\medterm{Opsoclonus} Rapid, involuntary, back-and-forth eye movements in many directions without pauses. It may occur with muscle jerks and poor coordination, sometimes because of an immune reaction, infection, medicine, or an underlying tumor, and needs medical evaluation.

\medterm{Opsonin} An immune protein that coats a germ and marks it for destruction. Antibodies and complement proteins can act as opsonins by helping infection-fighting cells recognize and engulf the germ.
\medterm{Opsonization} The coating of a germ or other particle with immune proteins so that infection-fighting cells can recognize, capture, and destroy it more easily. Antibodies and complement proteins commonly perform this function.

\medterm{Optic Atrophy} Damage and loss of optic-nerve fibers, causing permanent reduction in vision, color vision, or part of the visual field. Causes include glaucoma, poor blood supply, inflammation, pressure, injury, toxins, and inherited disorders; treatment targets the cause when possible.

\medterm{Optic Chiasm} An X-shaped junction beneath the brain where some fibers from the optic nerves cross. Damage can cause characteristic loss of parts of the visual field, depending on the location and cause.
\medterm{Optic Chiasma} The crossing point beneath the brain where some fibers from the two optic nerves switch sides. Damage there can cause loss of the outer parts of both visual fields.

\medterm{Optic Disc} The small area at the back of the eye where the optic nerve and blood vessels pass through the retina. It has no light-sensing cells, so it creates the eye’s normal blind spot.

\medterm{Optic Disc Edema} Swelling where the optic nerve enters the eye. It may result from raised pressure around the brain, inflammation, poor blood supply, or other eye or nerve disease. New swelling with headache or vision change requires urgent assessment.

\medterm{Optic Disk} The point where the optic nerve leaves the retina; it has no photoreceptors and forms the blind spot.

\medterm{Optic Glioma} A usually slow-growing tumor of the optic nerve or the visual pathway leading to the brain, most often in children. It may cause reduced vision, loss of part of the visual field, a bulging eye, or abnormal eye movements and is associated with neurofibromatosis type 1.

\synonyms
Optic Nerve Glioma
Optic Pathway Glioma

\medterm{Optic Nerve} The nerve that carries visual signals from each eye to the brain. Damage can cause reduced vision, loss of color or part of the visual field, and sometimes pain with eye movement.

\medterm{Optic Nerve Atrophy} Permanent damage and loss of nerve fibers carrying vision from the eye to the brain. It can reduce sharpness, color vision, or part of the visual field and may follow glaucoma, inflammation, poor blood supply, pressure, injury, or toxins.

\medterm{Optic Nerve Disorder} A disease affecting the optic nerve, which carries visual signals from the eye to the brain; it can cause reduced vision, color loss, or visual-field defects.

\medterm{Optic Nerve Glioma} Alternate index label; see \textit{Optic glioma}.

\medterm{Optic Neuritis} Inflammation or demyelination of the optic nerve, often causing subacute vision loss, impaired color vision, and pain with
eye movement. It may be associated with multiple sclerosis or other inflammatory,
infectious, or immune-mediated disorders.

\medterm{Optic Pathway Glioma} Alternate index label; see \textit{Optic glioma}.

\medterm{Optic Tract} A bundle of nerve fibers carrying visual information from the optic chiasma to brain areas that process sight and control visual reflexes. Damage causes loss of a characteristic part of the visual field.

\medterm{Optical Coherence Tomography} A painless scan that uses reflected light to create detailed cross-sectional pictures of the retina and other eye structures. It helps detect and monitor glaucoma, macular disease, swelling, and other changes without touching the eye.

\medterm{Optical Illusion} A visual perception in which an object, image, or movement appears different from its physical properties. Illusions are common, but new or persistent visual distortions may occur with eye, brain, medicine, or migraine-related disorders.
\medterm{Optical Rotation} The turning of plane-polarized light as it passes through a substance whose molecules have a particular three-dimensional arrangement. Measuring it can help identify or assess the concentration of some chemicals and medicines.

\medterm{Optician} A technician who fits and dispenses eyeglasses or contact lenses.

\medterm{Optometrist} A trained eye-care professional who examines vision and many ocular conditions and prescribes optical correction within local scope of practice.
\medterm{Optometry} The healthcare profession focused on examining vision and the eyes, prescribing glasses or contact lenses, detecting eye disease, and providing treatment or referral within the professional scope allowed by local law.

\medterm{Orthoptist} An eye-care professional who assesses and treats problems with eye movement, eye alignment, and the use of both eyes together. Orthoptists commonly help manage strabismus, amblyopia, double vision, and some vision problems in children and adults.

\medterm{Oral} Relating to the mouth, or taken through the mouth. Oral medicines are swallowed and absorbed through the digestive tract unless otherwise specified.
\medterm{Oral Administration} Taking a medicine by mouth so it is absorbed through the digestive tract; absorption can be affected by food, vomiting, interactions, or gastrointestinal disease.

\medterm{Oral Appliance Therapy} Treatment using a removable dental appliance during sleep, commonly to advance the mandible or retain the tongue and reduce upper-airway collapse in selected obstructive sleep apnea.

\medterm{Oral Cancer} Cancer of the lips, mouth, tongue, gums, inner cheeks, palate, or nearby throat. Tobacco, alcohol, betel nut, and some HPV infections increase risk; a persistent sore, lump, white or red patch, pain, or swallowing difficulty needs assessment.


\medterm{Oral Candidiasis} A yeast infection of the mouth, usually caused by Candida. It may produce wipeable white patches, soreness, burning, altered taste, or cracks at the mouth corners. Antibiotics, steroid inhalers, dentures, diabetes, dry mouth, and weak immunity increase risk.

\medterm{Oral Cavity} The space inside the mouth, extending from the lips to the throat and bounded by the cheeks, palate, tongue, and floor of the mouth. It receives food, begins digestion, and helps with speech, taste, and breathing.

\medterm{Oral Cavity Cancer} Cancer arising in the lips, tongue, gums, inner cheeks, floor of the mouth, or palate. Tobacco, alcohol, betel quid, and some HPV infections increase risk; persistent sores, lumps, patches, pain, or numbness need examination and often biopsy.

\medterm{Oral Cholecystogram} An older X-ray test in which swallowed contrast collects in the gallbladder, allowing gallstones and gallbladder function to be assessed. Ultrasound and other imaging methods have largely replaced it.

\medterm{Oral Contraceptive} A birth-control pill containing progestin alone or progestin combined with estrogen. It mainly prevents ovulation and thickens cervical mucus; it does not protect against sexually transmitted infections.

\medterm{Oral Glucose Tolerance Test} A diagnostic test measuring plasma glucose response after a standardized oral glucose
load. It is used in selected evaluations of diabetes, impaired glucose tolerance, and
gestational diabetes when fasting glucose or HbA1c is insufficient.

\medterm{Oral Health} Oral health is the condition of the teeth, gums, mouth, jaw, and related functions such as eating, speaking, and swallowing.

\medterm{Oral Hemorrhage} Bleeding from the mouth, gums, tongue, throat, or nearby tissues, caused by trauma, dental disease, infection, medicines, or a bleeding disorder.

\medterm{Oral Human Papillomavirus Infection} Oral human papillomavirus infection affects the mouth or throat; many infections clear without symptoms, while persistent high-risk types can contribute to some cancers.

\medterm{Oral Hygiene} Practices maintaining oral health: brushing twice daily with fluoride toothpaste,
flossing daily, mouthwash (optional), regular dental visits. Proper brushing technique:
soft brush, 45-degree angle to gumline, gentle circular motions. Interdental cleaning
essential for plaque removal.

\medterm{Oral Hypoglycemics Overdose} Taking too much diabetes medicine by mouth can cause dangerously low blood glucose, especially with insulin-releasing medicines. Sweating, shaking, confusion, seizures, or unconsciousness require immediate glucose treatment and emergency advice.

\medterm{Oral Leukoplakia} A persistent white patch in the mouth that cannot be rubbed off and is not explained by another condition. It is often linked to tobacco, alcohol, or betel-nut use and can sometimes develop into mouth cancer, so a persistent patch needs dental or medical assessment.

\medterm{Oral Lichen Planus} A long-lasting inflammatory condition causing lace-like white lines, red patches, or painful mouth sores, often on the cheeks or gums. It may fluctuate over time; persistent painful or changing areas need dental review to exclude other disease.

\synonyms
Lichen Planus, Oral
OLP

\medterm{Oral Moniliasis} Oral moniliasis is oral candidiasis, a Candida yeast infection causing white plaques, soreness, burning, or altered taste.


\medterm{Oral Mucosa} The moist lining of the mouth, including the inner lips and cheeks, gums, tongue, and palate. It protects deeper tissues, provides sensation, stays lubricated, and can develop ulcers, infections, inflammation, or tumors.


\medterm{Oral Mucous Cyst} An oral mucous cyst is a mucus-filled lesion caused by salivary-gland duct injury or blockage, often appearing on the lower lip or floor of the mouth.

\medterm{Oral Nutritional Supplement} A calorie-, protein-, or micronutrient-containing product consumed by mouth to
supplement, not automatically replace, usual food intake. The product and dose should
match the patient's needs, swallowing ability, disease restrictions, and tolerance.

\medterm{Oral Pathology} The study and diagnosis of diseases affecting the mouth, teeth, jaws, and related tissues. These include tooth decay, infections, cysts, inflammation, developmental changes, and tumors. Examination, imaging, laboratory tests, or a tissue sample may be needed to identify the cause.

\medterm{Oral Potentially Malignant Disorders} Conditions affecting the mouth lining that carry an increased risk of developing into oral cancer. Examples include leukoplakia, erythroplakia, and some long-lasting inflammatory or inherited disorders; appearance alone cannot determine the risk.
\medterm{Oral Radiology} The use and interpretation of X-rays and other scans of the mouth, teeth, jaws, face, and nearby structures. These images help find decay, infection, fractures, tumors, bone disease, and problems related to treatment planning.

\medterm{Oral Rehydration Salts} A balanced mixture of glucose and electrolytes dissolved in clean water to replace fluid and salts lost through diarrhea or vomiting.

\medterm{Oral Rehydration Therapy} Replacement of water and electrolytes by mouth using a balanced glucose-and-salt solution, especially to prevent or treat dehydration from diarrhoea or vomiting.

\synonyms
ORT

\medterm{Oral Sex (Fellatio, Cunnilingus)} Oral stimulation of the genitals for sexual pleasure. Fellatio: oral stimulation of the
penis. Cunnilingus: oral stimulation of the vulva and clitoris. Analingus (rimming):
oral stimulation of the anus.

\synonyms
Fellatio, Cunnilingus

\medterm{Oral Spray} A liquid medicine or preparation sprayed into the mouth or throat for local treatment or absorption through oral tissues.

\medterm{Oral Squamous Cell Carcinoma} The most common type of mouth cancer, arising from the cells lining the lips, tongue, gums, cheeks, or palate. A sore, lump, red or white patch, bleeding, pain, or numbness that persists needs examination and often a biopsy.

\medterm{Oral Surgery} Operations on the mouth, teeth, jaws, or nearby facial tissues. Procedures include tooth removal, dental implants, tissue sampling, cyst removal, repair of injury, and correction of jaw problems.

\medterm{Oral Thrush} A fungal infection of the mouth, usually caused by Candida, producing creamy white patches, soreness, or altered taste.

\synonyms
Candidiasis, Oral
Thrush, Oral

\medterm{Oral Toxicity} Harm to the mouth, throat, or digestive system caused by a swallowed poison, medicine, treatment, or excessive dose. It may cause burning, sores, nausea, vomiting, abdominal pain, or organ injury and may require urgent poison-control or medical advice.

\medterm{Orange Peel Skin} Skin with small dimples resembling an orange peel because swelling stretches the skin around hair follicles. It may result from blocked lymph drainage, especially in breast disease, and a new breast change needs assessment.
\medterm{Orbicularis} A circular muscle surrounding an opening, such as orbicularis oculi around the eye or orbicularis oris around the mouth.
\medterm{Orbit} The bony socket that contains and protects the eye, its muscles, nerves, blood vessels, and supporting tissues. Problems in the orbit can cause pain, swelling, double vision, bulging, or impaired eye movement.
\medterm{Orbit CT Scan} A CT scan of the eye socket that uses X-rays and computer images to show the bones, eye muscles, nerves, blood vessels, and surrounding tissues. It can help assess injury, infection, bleeding, or a growth.

\medterm{Orbital} Relating to the orbit, the bony cavity containing the eye, or to structures within it. Orbital disease may affect the eye, muscles, nerves, or surrounding tissues.
\medterm{Orbital Cellulitis} A serious infection in the tissues around and behind the eye. It can cause eyelid swelling, fever, pain with eye movement, double vision, or reduced vision and can spread to the brain, so it requires urgent hospital treatment.

\medterm{Orbital Contusion} A bruise of the tissues around the eye caused by blunt trauma, often producing swelling, bruising, and tenderness.

\medterm{Orbital Fracture} A break in one or more bones forming the eye socket, usually caused by facial trauma and potentially associated with double vision, eye-movement restriction, or injury to the eye.

\medterm{Orbital Fractures} Breaks in the bones surrounding the eye, which may cause swelling, double vision, numbness, restricted eye movement, or injury to the eye itself.

\medterm{Orbital Irradiation} Radiation treatment directed at tissues within the eye socket, sometimes used for severe thyroid eye disease.

\medterm{Orbital Metastasis} Cancer that has spread to the bony socket around the eye from another organ, commonly the breast, lung, or prostate. It may cause bulging of the eye, pain, double vision, restricted movement, or reduced sight.

\medterm{Orbital Myositis} Inflammation of one or more muscles that move the eye. It can cause pain, eyelid swelling, double vision, and restricted eye movement; treatment depends on the cause and may require urgent specialist assessment.

\medterm{Orbital Neoplasm} An abnormal growth in the eye socket that may arise from the eye, optic nerve, muscles, glands, bone, or other tissues. It can cause bulging, pain, double vision, or vision loss and may be benign or cancerous.

\medterm{Orbital Pseudotumor} Painful inflammation of tissues around the eye that is not caused by infection or a true tumor. It may cause eyelid swelling, bulging of the eye, double vision, and painful or limited eye movement.

\synonyms
idiopathic.


\medterm{Orchialgia} Pain in one or both testicles. Sudden severe pain, swelling, nausea, or a high-riding testicle may indicate twisting of the testicle and requires emergency assessment. Longer-lasting pain can result from infection, injury, varicocele, nerve pain, or other causes.
\medterm{Orchidectomy} Surgical removal of one or both testes, also called orchiectomy.
\medterm{Orchiectomy} Surgical removal of one or both testicles. Radical inguinal orchiectomy is diagnostic
and therapeutic for testicular cancer; bilateral orchiectomy is a hormonal therapy for
prostate cancer.

\medterm{Orchitis} Orchitis is inflammation of one or both testes, often from infection such as mumps or bacteria, causing pain, swelling, and sometimes impaired fertility.

\medterm{Order Coleoptera} The insect order containing beetles, characterized by hardened forewings that protect the flying wings; some species affect health as pests or vectors.

\medterm{Order Spirochaetales} A biological order of spiral-shaped bacteria that includes groups causing syphilis, Lyme disease, relapsing fever, and leptospirosis. Identification of the species determines medical significance and treatment.
\medterm{Orexin} Orexin, or hypocretin, is a hypothalamic neuropeptide that promotes wakefulness, arousal, appetite, and energy regulation; loss of orexin signaling is strongly associated with narcolepsy type 1.

\medterm{Orexin Receptor Antagonists} Medications that block orexin (hypocretin) receptors to promote sleep. Suvorexant and
lemborexant are approved for insomnia. They reduce wakefulness by blocking the arousal
system. They have lower abuse potential than benzodiazepines and do not cause
respiratory depression.

\medterm{Orf} A zoonotic viral skin infection caused by the parapoxvirus, transmitted to humans from
infected sheep and goats (direct contact with lesions, fomites—virus survives in scabs
for months). Orf is an occupational disease of farmers, butchers, and veterinarians.
Incubation: 3--7 days.

\medterm{Organ} A structure composed of two or more different types of tissues that has a recognizable
shape and performs specific functions.


\medterm{Organ Failure} Severe loss of an organ's ability to perform essential functions, such as filtering blood, exchanging oxygen, or pumping blood.

\medterm{Organ System} A group of organs that work together to perform an important body function, such as breathing or digestion.

\medterm{Organ Transplant} Surgical transfer of an organ from a donor to a recipient to replace failing organ function.
\medterm{Organ Transplantation} Replacement of a failing organ with a healthy organ from a donor. Transplants can restore function but require careful matching, lifelong monitoring, and medicines to prevent rejection; infection, rejection, and medicine-related harm remain possible.

\medterm{Organ Transplants} The surgical transfer of an organ from a donor to a recipient whose organ function is severely
impaired. Solid-organ transplantation may involve the kidney, liver, heart, lung, pancreas,
or intestine and requires donor assessment, recipient selection, and ongoing follow-up.

\medterm{Organelle} A specialized part inside a cell that performs a particular job. Examples include the nucleus, which stores genetic instructions, and mitochondria, which release energy; organelles work together to keep the cell alive.
\medterm{Organic} Relating to living matter or carbon-containing chemistry; in psychiatry, an older term for symptoms caused by a physical brain disorder.

\medterm{Organic Chemical} A carbon-containing chemical compound, usually also containing hydrogen and sometimes oxygen, nitrogen, sulfur, or other elements. Organic chemicals include medicines, fats, sugars, proteins, solvents, and many toxins.

\medterm{Organic Disease} An illness caused by a physical change or measurable disturbance in the body, such as infection, inflammation, tissue damage, or abnormal organ function. The term is older and is less useful when symptoms have not yet been explained.
\medterm{Organic Impurity} An organic impurity is an unintended carbon-containing chemical in a medicine, material, or specimen that may affect safety, quality, or measurement.


\medterm{Organization Name} The formal name of a group, institution, company, or agency; in healthcare records it should identify the responsible organization accurately.

\medterm{Organogenesis} The process during embryonic development in which groups of cells form and organize into organs. Disruption can lead to congenital abnormalities, while the process is also studied in tissue engineering and regenerative medicine.



\medterm{Organophosphate} A class of chemicals that can inhibit acetylcholinesterase. Some are pesticides or nerve agents and can cause cholinergic poisoning.

\medterm{Organophosphate Poisoning} A serious poisoning caused by pesticides or nerve agents that excessively stimulate nerves and muscles. It may cause sweating, drooling, vomiting, diarrhea, pinpoint pupils, muscle twitching, wheezing, weakness, seizures, or breathing failure. Remove contaminated clothing, wash the skin, and seek emergency care immediately.

\medterm{Organosilicon Compounds} Organosilicon compounds contain silicon bonded to carbon and are used in materials, medicines, devices, and industry; biological effects depend on the specific compound.

\medterm{Organotherapy} An outdated treatment method using extracts or tissues from organs, with little accepted modern medical use.
\medterm{Orgasm} The peak of sexual pleasure, usually accompanied by involuntary muscle and body responses that vary among individuals.
\medterm{Orgasm Male (Sexual (Sex) Problems In Men)} Alternate index label for Sexual (Sex) Problems in Men.

\medterm{Orgasmic Dysfunction} Persistent or recurrent difficulty reaching orgasm, markedly reduced orgasmic intensity, or absent orgasm that causes clinically significant distress. Medical, psychological, medication-related, hormonal, and relationship factors may contribute.
\medterm{Orgasmic Dysfunction In Women} Alternate index label for Orgasmic Dysfunction.

\medterm{Orientation} Awareness of personal identity, place, time, and situation. Clinicians assess orientation as part of mental-status and neurologic examinations.
\medterm{Orifice} An opening leading into or out of a body structure, organ, or passage, allowing materials to enter or leave.

\medterm{Orlistat} A medicine used with a reduced-calorie diet to help with weight loss or weight control. It stops some dietary fat from being absorbed, so oily stools, gas, and urgent bowel movements can occur; it can also reduce absorption of vitamins and some medicines.
\medterm{Ornithine} A nonessential amino acid and intermediate in the urea cycle that helps convert toxic ammonia into urea.

\medterm{Ornithine Decarboxylase} An enzyme that converts ornithine to putrescine, an early step in polyamine synthesis involved in cell growth.

\medterm{Ornithine Transcarbamylase Deficiency} OTC deficiency is an inherited urea-cycle disorder causing ammonia accumulation, vomiting, confusion, seizures, or coma, especially during metabolic stress.

\medterm{Orogastric Tube} A flexible tube passed through the mouth into the stomach. It may deliver nutrition or medicines, remove stomach contents, or decompress the stomach when nasal placement is unsuitable.

\medterm{Oropharyngeal} Oropharyngeal means relating to the mouth and the part of the throat behind it, including structures involved in breathing, swallowing, and speech.

\medterm{Oropharyngeal Airway} A rigid device placed through the mouth to keep the tongue from blocking the throat in an unconscious person. It must not be used in someone with a gag reflex and does not prevent vomit from entering the lungs.

\medterm{Oropharyngeal Carcinoma} Cancer arising in the part of the throat behind the mouth, including the tonsils, back of the tongue, soft palate, or throat wall. Tobacco, alcohol, and some HPV infections increase risk; assessment usually includes examination, tissue sampling, and scans.

\medterm{Oropharyngeal Dysphagia} Difficulty moving food or liquid from the mouth through the throat. It may cause coughing, choking, a wet voice, food coming through the nose, or food entering the airway and needs assessment when persistent.

\medterm{Oropharyngeal Neoplasm} An oropharyngeal neoplasm is an abnormal growth in the throat behind the mouth, including the tonsils and base of the tongue; it may be benign or malignant.

\medterm{Oropharynx} The part of the throat behind the mouth, including the tonsils and back of the tongue. It helps move food and liquid toward the esophagus and carries air toward the voice box.

\medterm{Oropharynx Lesion Biopsy} Sampling of an oropharyngeal lesion, usually through the mouth under local or procedural anesthesia, for histopathologic evaluation of infection, dysplasia, or cancer. Persistent or suspicious lesions require specialist assessment.

\medterm{Orotic Aciduria} An unusually high amount of orotic acid in urine. It may result from a rare inherited disorder affecting blood-cell and DNA building blocks, causing anemia and poor growth, or from a urea-cycle disorder with high ammonia.
\medterm{Orphan Disease} A rare disease affecting relatively few people; definitions and eligibility for orphan-drug programs vary by jurisdiction.

\medterm{Orphenadrine} A centrally acting medicine with anticholinergic properties used in some countries for painful muscle spasm; it can cause drowsiness, dry mouth, and confusion.



\medterm{Orthodontic Appliance} A device used to move teeth, guide jaw growth, maintain alignment, or correct a bite problem.

\medterm{Orthodontic Bracket} A bonded or banded attachment that transmits forces from an orthodontic archwire to a
tooth. Bracket position and wire sequence guide controlled tooth movement.

\medterm{Orthodontic Retainer} A removable or fixed appliance used after active orthodontic treatment to maintain tooth
position while periodontal and alveolar tissues remodel. Retention is often needed long
term because relapse can occur.

\medterm{Orthodontic Treatment (Dental Braces)} Treatment that uses braces, wires, aligners, or other appliances to move teeth and guide jaw development. It can improve bite, spacing, and cleaning access; treatment requires assessment and long-term retention.

\medterm{Orthodontic Wires} Metal or other wires used with braces to apply controlled forces that move teeth or help maintain their position.

\medterm{Orthodontics} The dental specialty that corrects teeth or jaws that do not align properly. Braces, clear aligners, or other appliances can improve biting, chewing, appearance, and cleaning; retainers are usually needed afterward because teeth can shift back.

\medterm{Orthodontist} A dentist specially trained to diagnose and treat teeth or jaws that do not align properly. Orthodontists use braces, clear aligners, retainers, and sometimes surgery to improve biting, chewing, appearance, and oral health.
\medterm{Orthognathic Surgery} Corrective surgery to reposition the upper jaw, lower jaw, or both when bone alignment causes problems with biting, chewing, speaking, breathing, or facial development. It is often combined with orthodontic treatment.

\medterm{Orthohepadnavirus} A group of viruses that includes hepatitis B virus and related viruses infecting mammals. These viruses carry genetic material as DNA and may cause long-lasting liver infection, inflammation, scarring, or cancer.

\medterm{Orthokeratosis} Thickening of the outer skin layer with fully mature skin cells. It can be a normal protective response, such as on the soles, or part of conditions causing thick, scaly, or rough skin.

\medterm{Orthomyxoviridae} A family of viruses that includes influenza A, B, and C. They mainly infect the airways and can cause flu, outbreaks, and sometimes severe pneumonia or other complications.
\medterm{Orthopaedic Surgery} Surgery for disorders and injuries involving bones, joints, muscles, ligaments, tendons, and related tissues. Procedures may repair fractures, replace joints, correct deformity, stabilize the spine, or relieve pressure on nerves.
\medterm{Orthopaedics} The British spelling of orthopedics, the medical and surgical specialty concerned with the musculoskeletal system.
\medterm{Orthopedic Biologics} Biological materials used to help bones, tendons, or other tissues heal. They include a person’s own bone, donated bone, processed tissue, and synthetic substitutes; benefits and risks vary by the injury and procedure.

\medterm{Orthopedic Implant} A device placed in or on the body to repair, support, replace, or stabilize a bone or joint. Examples include plates, screws, rods, joint replacements, and artificial ligaments. Possible complications include infection, loosening, breakage, wear, or bone damage.

\medterm{Orthopedic Infections} Infections involving bones, joints, muscles, or orthopedic implants. Examples include osteomyelitis, septic arthritis,
prosthetic-joint infection, and some diabetic foot infections.

\medterm{Orthopedic Rehabilitation} Treatment that restores movement, strength, daily function, and independence after injury or surgery involving bones, joints, muscles, ligaments, or tendons. It may include exercises, activity training, pain management, and assistive equipment.
\medterm{Orthopedic Screws} Metal or other screws used to hold broken bones or surgical implants in position while bone heals. Their shape and size are chosen for the bone and purpose. Complications can include loosening, breakage, infection, or injury to nearby structures.

\medterm{Orthopedic Services} Orthopedic services diagnose and treat problems of bones, joints, muscles, ligaments, tendons, and related movement systems.

\medterm{Orthopedics} The branch of surgery concerned with the musculoskeletal system, including bones,
joints, ligaments, tendons, muscles, and nerves.

\medterm{Orthopedist} A physician who diagnoses and treats problems of bones, joints, muscles, ligaments, tendons, and related tissues. Treatment may include rehabilitation, medicines, injections, or surgery.
\synonyms
Orthopedic surgeon

\medterm{Orthopnea} Shortness of breath that begins or worsens when lying flat and improves when sitting or standing. It often occurs with heart or lung disease, but other causes are possible; new or worsening symptoms need medical assessment.

\medterm{Orthopneic Position} Sitting upright, often leaning forward with the arms supported, to make breathing easier. A person who suddenly needs this position or cannot lie flat because of breathlessness should seek medical assessment.

\medterm{Orthopnoea} Shortness of breath that starts or worsens when lying flat and improves when sitting or standing. It often occurs with heart or lung disease; new or worsening difficulty lying down because of breathlessness needs medical assessment.


\medterm{Orthopoxvirus} A genus of poxviruses that includes variola, mpox, vaccinia, and cowpox viruses. Infection may cause fever and skin lesions; severity and transmission depend on the specific virus.
\medterm{Orthoptera} The insect order containing grasshoppers, crickets, katydids, and locusts; some species cause bites, allergies, or crop-related health effects.

\medterm{Orthoptic} Relating to assessment and treatment of binocular vision, eye alignment, and disorders such as strabismus.
\medterm{Orthoreovirus} A genus of nonenveloped, segmented double-stranded RNA viruses; human infection is often mild or asymptomatic but can occasionally cause disease.

\medterm{Orthorexia} A term used for an excessive or rigid preoccupation with eating foods considered healthy. It is not a universally recognized formal psychiatric diagnosis but may cause nutritional or social impairment.

\medterm{Orthoscopic} Orthoscopic means providing a straight or undistorted view, as in an optical instrument or examination.
\medterm{Orthosis} A brace, splint, shoe insert, or similar device worn outside the body to support, align, protect, or improve movement of a body part. Unlike a prosthesis, it does not replace a missing body part.

\medterm{Orthostatic Hypotension} A fall in blood pressure after standing, often causing dizziness, blurred vision, weakness, or fainting. It may result from dehydration, medicines, autonomic nerve problems, or prolonged bed rest. Repeated symptoms or falls need medical assessment.

\medterm{Orthostatic Intolerance} Dizziness, fast heartbeat, weakness, blurred vision, fatigue, or faintness that develops or worsens on standing because the body does not maintain blood flow to the brain adequately. Causes include dehydration, medicines, nerve disorders, and some chronic conditions.

\medterm{Orthotic} Orthotic means relating to a brace, splint, insert, or device that supports, aligns, protects, or improves function of a body part.

\medterm{Orthotic Devices} External supports that hold, align, protect, or assist a body part, such as a brace, splint, shoe insert, or spinal support. They can reduce pain and improve movement, but poor fit may cause pressure, numbness, skin injury, or restricted circulation.

\medterm{Orthotic Training} Instruction and practice for safely using a brace or other support. It includes putting it on and removing it, checking skin, cleaning it, walking with it, and recognizing problems needing review.

\medterm{Ortolani's Sign} A palpable or audible clunk produced when a dislocated but reducible infant hip is gently abducted and reduced; it is a screening sign of developmental dysplasia of the hip.

\medterm{Ortolani's Test} A hip examination maneuver in infants that detects a reducible dislocation by producing a palpable clunk as the femoral head enters the acetabulum.
\medterm{Os} A bone or bony structure; in Latin anatomical terminology it means bone, while in obstetrics it means an opening.
\medterm{Os Odontoideum} Os odontoideum is a separate ossicle near the second cervical vertebra’s odontoid process that can cause atlantoaxial instability or spinal-cord risk.

\medterm{Osborn Wave} A notch or extra wave on an electrocardiogram, classically seen when the body is dangerously cold. It can also occur in other conditions and must be interpreted with the person's temperature and other findings.

\synonyms
J wave

\medterm{Oscillation} Repeated movement or variation around a central position or value.
\medterm{Oscillator} A device or system that produces repeated movement or an electrical signal at regular intervals.
\medterm{Oscillograph} An instrument or computer display that records a changing signal over time. It can show electrical, sound, pressure, or movement waves and has been used in medical tests to study signals from the heart, brain, muscles, or breathing.
\medterm{Oscillometry} A method that measures pressure-related vibrations, commonly used by automated blood-pressure cuffs and some vascular devices. The result estimates blood pressure or blood-flow characteristics but may be inaccurate with movement or irregular rhythms.
\medterm{Oscillopsia} The disturbing illusion that the visual environment is moving or bouncing, often caused by vestibular or eye-movement dysfunction.
\medterm{Oscitation} Yawning, a reflex involving opening the mouth and deep inhalation that is associated with sleepiness, fatigue, or other physiologic states.

\medterm{Oseltamivir} Oseltamivir is an antiviral medicine used to treat or prevent influenza, with the greatest benefit when started early.

\medterm{Oseltamivir Phosphate} Oseltamivir phosphate is the salt form of oseltamivir, an antiviral medicine used for treatment or prevention of influenza.

\medterm{OSFED} A recognized eating disorder in which symptoms cause significant distress or interfere with health or daily life but do not meet every criterion for another named eating disorder. It still requires assessment and treatment.

\medterm{Osgood-Schlatter Disease} Painful irritation of the growing area where the kneecap tendon attaches to the shinbone. It commonly affects active children and teenagers and causes tenderness or swelling below the knee.

\medterm{Osler Nodes} Tender red or purple lumps on the pads of the fingers or toes. They can occur with infective endocarditis and some other conditions; they are an important clue but do not prove the diagnosis, which requires a full medical assessment.

\medterm{Osmidrosis} Strong, unpleasant body odor caused when skin bacteria break down sweat and skin oils. It commonly affects the armpits or feet and may improve with washing, dry clothing, antiperspirants, treatment of skin infection, or assessment of an underlying condition.
\medterm{Osmium} A dense metallic element; some osmium compounds are highly toxic and can damage the eyes and respiratory tract.
\medterm{Osmolality} A measure of how many dissolved particles are present in a kilogram of fluid. Blood osmolality depends mainly on sodium, glucose, and urea and helps determine how water moves between the bloodstream and cells.

\medterm{Osmolality Blood Test} A blood osmolality test measures the concentration of dissolved particles in blood and helps assess hydration, electrolyte disorders, and certain poisonings.

\medterm{Osmolality Urine Test} A urine osmolality test measures how concentrated the urine is and helps evaluate the kidneys' ability to regulate water balance.

\medterm{Osmolar} Describing a liquid according to the number of dissolved particles in each liter. This property affects water movement between body compartments and helps clinicians assess blood chemistry, hydration, kidney function, and intravenous fluids.

\medterm{Osmolarity} The concentration of osmotically active particles per liter of solution.
\medterm{Osmole} A unit representing the amount of a substance that contributes one mole of osmotically active particles.

\medterm{Osmoreceptor} A sensory cell or neural receptor that detects changes in osmotic concentration, especially in the hypothalamus.
\medterm{Osmoregulation} Control of water and dissolved-solute balance in the body or an organism, maintaining appropriate cell volume and internal concentration.

\medterm{Osmosis} Movement of water through a membrane toward the side with more dissolved particles. This process helps determine how water is distributed between blood, cells, and surrounding tissues.

\medterm{Osmotic Demyelination Syndrome} Serious brain injury that can occur when long-standing low blood sodium is corrected too quickly. Symptoms may appear days later and include confusion, difficulty speaking or swallowing, weakness, movement problems, or reduced consciousness.

\medterm{Osmotic Diarrhea} Diarrhea caused by poorly absorbed solutes retaining water in the intestinal lumen. It
typically improves with fasting and may result from carbohydrate malabsorption,
laxatives, or other osmotically active substances.

\medterm{Osmotic Diuresis} Increased urine production caused by poorly reabsorbed solutes in the kidney tubules, such as glucose during marked hyperglycemia.

\medterm{Osmotic Fragility Test} A test that measures the susceptibility of red blood cells to haemolysis when placed in
progressively more dilute saline solutions. Increased fragility supports hereditary
spherocytosis, although the result may also be affected by other haemolytic disorders
and recent transfusion.

\medterm{Osmotic Pressure} The force that draws water across a membrane toward the side containing more dissolved particles. It helps control how water moves between blood, cells, and surrounding tissues.

\medterm{Osseointegration} The firm attachment of living bone to the surface of an implant, such as a dental implant or artificial joint. Good stability, healthy bone, healing, and infection control support this process.

\medterm{Osseous} Relating to or composed of bone. For example, an osseous lesion is a bone lesion.

\medterm{Osseous Metaplasia} Osseous metaplasia is formation of mature bone in tissue where bone is not normally present, sometimes after injury, inflammation, or within a lesion.

\medterm{Ossicle} A small bone, especially one of the three tiny middle-ear bones. The malleus, incus, and stapes transmit eardrum vibrations to the inner ear, allowing sound to be heard.
\medterm{Ossicles} Three smallest body bones forming the ossicular chain in the middle ear: malleus
(hammer, attached to tympanic membrane), incus (anvil), stapes (stirrup, footplate in
oval window).

\medterm{Ossicular Discontinuity} Partial or complete separation of the middle-ear ossicles, usually caused by trauma, chronic ear disease, or surgery.
It produces conductive hearing loss and may be evaluated with otoscopy, audiometry, and imaging.

\synonyms
Ossicular-chain discontinuity; ossicular disruption.

\medterm{Ossicular Fusion} Abnormal joining or fixation of the middle-ear ossicles (malleus, incus, or stapes), which can restrict sound transmission and cause conductive hearing loss. Evaluation and treatment depend on the underlying congenital, inflammatory, traumatic, or postoperative cause.
\medterm{Ossification} The formation of bone. It occurs during growth and fracture healing and can also happen abnormally in soft tissues after severe injury, burns, or nervous-system damage.
\medterm{Ossifying Fibroma} An ossifying fibroma is a usually benign fibro-osseous tumor that replaces normal bone with fibrous tissue and mineralized material, often in the jaw.

\medterm{Osteitis} Inflammation of bone caused by infection, injury, inflammatory disease, or another process. It may cause localized pain, swelling, warmth, reduced movement, or fever and sometimes requires imaging, antibiotics, or surgery.
\medterm{Osteitis Deformans} Osteitis deformans is Paget disease of bone, in which abnormal remodeling produces enlarged, weakened, or deformed bones.

\medterm{Osteitis Deformans (Paget's Disease)} Alternate index label for Paget's Disease.

\medterm{Osteitis Fibrosa} Osteitis fibrosa is fibrous replacement and weakening of bone caused by prolonged excess parathyroid hormone, usually in severe hyperparathyroidism.

\medterm{Osteoarthritis} A joint disease in which cartilage and other tissues gradually break down or change. It can cause pain, stiffness, swelling, bone changes, and restricted movement. It usually develops over time but may occur earlier after joint injury, abnormal joint structure, or previous joint disease.

\synonyms
Arthritis Degenerative (Osteoarthritis)
Degenerative Arthritis

\medterm{Osteoarthritis (Geriatric)} Osteoarthritis occurring in later life, commonly affecting the knees, hips, hands, or spine
and causing joint pain, stiffness, and reduced movement. Management may include exercise,
weight management, medicines, assistive measures, or joint replacement when appropriate.

\synonyms
Osteoarthritis in older adults; age-related osteoarthritis

\medterm{Osteoarthritis, Cervical} Alternate index label; see \textit{Cervical spondylosis}.

\medterm{Osteoarthropathy} Osteoarthropathy is disease affecting bones and joints, often describing degenerative, hypertrophic, metabolic, or systemic-associated changes.
\medterm{Osteoarthrosis} Another name for osteoarthritis, a joint disorder in which cartilage and nearby bone gradually change, causing pain, stiffness, and reduced movement.
\medterm{Osteoblast} A cell that builds new bone. It makes the soft framework of bone and helps add calcium and phosphate to harden it. Some osteoblasts become osteocytes, which maintain mature bone.

\medterm{Osteoblastoma} A rare benign bone-forming tumor, usually in the spine or long bones, that can cause persistent pain and sometimes neurologic symptoms.

\medterm{Osteoblasts} Cells that build new bone. They produce the material that forms bone and help control its hardening with calcium and phosphate. Some osteoblasts become osteocytes, which help maintain mature bone.

\medterm{Osteocalcin} A protein made mainly by osteoblasts that participates in bone mineralization and serves as a marker of bone formation.

\medterm{Osteocalcin Assay} A blood test measuring osteocalcin, a protein made by bone-forming cells. It can provide information about bone formation and turnover, but results are affected by age, kidney function, medicines, and the laboratory method.
\medterm{Osteochondritis} Inflammation or injury involving bone and its overlying cartilage, often near a joint. The term includes conditions such as osteochondritis dissecans, in which a fragment may loosen.
\medterm{Osteochondritis Dissecans} An area of damaged bone and cartilage beneath a joint surface that may loosen or separate. It most often affects the knee in adolescents and young adults and can cause activity-related pain, swelling, catching, or locking. Treatment depends on lesion stability, age, and symptoms.

\medterm{Osteochondroma} A usually noncancerous cartilage-covered growth projecting from the surface of a bone, often near a growth plate. It may be single or occur as multiple inherited growths and can cause pain, deformity, or pressure on nearby nerves or vessels.

\medterm{Osteochondromatosis} Osteochondromatosis is formation of multiple cartilage-capped bony outgrowths, usually around growth plates, that can cause deformity, pain, or nerve and vessel compression.

\medterm{Osteochondrosis} A disorder of bone and cartilage development, often involving a growth area or joint surface. It may cause pain, swelling, stiffness, or altered movement and can affect children, adolescents, or adults.
\medterm{Osteoclasis} Deliberate surgical breaking or manipulation of a bone to correct a deformity or restore alignment. It requires stabilization while the bone heals and may cause bleeding, infection, or nerve injury.
\medterm{Osteoclast} A specialized bone cell that breaks down old bone during normal renewal and repair. Osteoclast activity helps regulate bone strength and blood calcium; excessive activity can contribute to osteoporosis and other diseases causing bone loss.
\medterm{Osteoclasts} Cells that break down old bone. This normal process helps bones grow, repair themselves, and release calcium into the blood. Too much osteoclast activity causes bone loss, as in osteoporosis, Paget disease, and some cancers that spread to bone. Medicines such as bisphosphonates and denosumab can slow osteoclast activity and reduce bone weakening.

\medterm{Osteocyte} A mature bone cell embedded in the hard bone it helped form. Osteocytes sense strain and chemical changes and send signals that help control bone strength, repair, and removal of old bone.

\medterm{Osteocytes} Mature bone cells trapped inside the hard bone they helped form. They sense strain and changes in the bone and send signals that help control bone repair, strengthening, and breakdown.

\medterm{Osteodystrophy} Abnormal bone growth, strength, or renewal caused by a body-wide or metabolic disorder. It may result from kidney disease, mineral imbalance, hormone problems, or inherited conditions and can cause pain or fractures.

\medterm{Osteogenesis} The formation and development of bone. It occurs in the growing skeleton and during fracture repair. Abnormal osteogenesis can contribute to birth disorders, poor healing, or fragile bones.
\medterm{Osteogenesis Imperfecta} A group of genetic disorders, also called brittle bone disease, in which bones break easily, sometimes after mild injury or for no clear reason. The disorders are often caused by gene changes that reduce or alter type I collagen, a protein that gives strength to bone and other connective tissues. Osteogenesis imperfecta is present from birth and can range from mild to severe; severe forms may cause fractures before birth. Other features may include short height, loose joints, blue or gray whites of the eyes, hearing loss, weak or discolored teeth, weak muscles, or a curved spine.

\medterm{Osteogenesis Imperfecta Brittle Bone Disease} Alternate index label; see \textit{Osteogenesis Imperfecta}.

\synonyms
Brittle Bone Disease (Osteogenesis Imperfecta Brittle Bone Disease)

\medterm{Osteogenic Sarcoma} A cancerous bone tumor whose cells make immature bone tissue. It often develops near the ends of growing long bones and may cause persistent pain, swelling, reduced movement, or a fracture.

\medterm{Osteoid} The newly formed, unmineralized organic framework of bone, produced mainly by osteoblasts. Minerals are later deposited in it to make mature hard bone.
\medterm{Osteoid Osteoma} A small noncancerous bone tumor, usually in an adolescent or young adult. It often causes persistent pain that is worse at night and may improve with anti-inflammatory medicine; scans help locate it and guide treatment.

\medterm{Osteology} The study of bones, including their structure, development, function, injuries, and diseases. It supports the assessment and treatment of fractures, bone deformities, weak bones, infections, and bone tumors.
\medterm{Osteolysis} Osteolysis is progressive destruction or resorption of bone caused by tumors, inflammation, prosthesis reactions, infection, metabolic disease, or other processes.

\medterm{Osteolytic} Causing or characterized by breakdown and loss of bone tissue. Osteolytic lesions can weaken bones and occur with cancers, infections, inflammation, metabolic disorders, or excessive activity of bone-removing cells.

\medterm{Osteoma} A benign, slow-growing tumor made of bone. It often develops in the skull or facial bones and may cause no symptoms unless it blocks a passage or presses on nearby structures.
\medterm{Osteoma Cutis} Osteoma cutis is formation of bone within the skin or subcutaneous tissue, occurring as a primary disorder or secondary to acne, trauma, or another condition.

\medterm{Osteomalacia} Osteomalacia is inadequate mineralization of adult bone, commonly from vitamin-D deficiency or phosphate disorders, causing bone pain, weakness, and fractures.

\medterm{Osteomalacia (Rickets)} Alternate index label for Rickets.

\medterm{Osteomyelitis} Infection and inflammation of bone, usually caused by bacteria and sometimes fungi or other organisms. It can cause pain, fever, swelling, reduced movement, bone damage, or a persistent draining wound.


\medterm{Osteomyelitis In Children} Childhood osteomyelitis is infection of bone, often caused by bacteria entering through blood or trauma, producing pain, fever, swelling, and impaired movement.

\medterm{Osteon} A tube-shaped structural unit of hard, compact bone arranged around a small channel containing blood vessels and nerves. Osteons help give bone strength while allowing it to receive nutrients.

\medterm{Osteonecrosis} Osteonecrosis is death of bone tissue from inadequate blood supply, causing pain, collapse, and joint damage, commonly in the hip or knee.

\medterm{Osteonecrosis (Avascular Necrosis)} Death of bone tissue after its blood supply is reduced, most often affecting the hip, knee, shoulder, or ankle. It can cause pain and collapse of the joint and may follow injury, prolonged steroid use, heavy alcohol use, blood disorders, or other conditions.

\synonyms
Avascular Necrosis

\medterm{Osteonecrosis Of The Jaw} Death of jawbone tissue, often associated with antiresorptive or antiangiogenic medications, radiation, infection, or trauma, and characterized by exposed bone or a nonhealing oral wound.

\medterm{Osteonectin} An extracellular-matrix protein involved in collagen organization, mineral binding, and tissue remodeling; it is also called SPARC.

\medterm{Osteons} Cylindrical structural units of compact bone arranged around central canals.

\synonyms
Haversian systems

\medterm{Osteopath} A physician trained in medicine and surgery who earns a Doctor of Osteopathic Medicine degree.
\synonyms
DO

\medterm{Osteopathy} A term used differently in different countries. It may refer to physician-led osteopathic medicine or to manual treatment focused on the muscles and joints; the practitioner’s training and legal scope should be checked locally.

\medterm{Osteopenia} Bone density that is lower than average but not low enough to meet the definition of osteoporosis. It may increase fracture risk, especially with age, previous fractures, falls, medicines, or other conditions that weaken bone.

\medterm{Osteopenia - Premature Infants} Osteopenia of prematurity is reduced bone mineralization in premature infants due to immature mineral stores, limited intake, illness, or prolonged intensive care.

\medterm{Osteopetrosis} A rare inherited disorder in which old bone is not removed normally. Bones become unusually dense but brittle and may fracture; crowding of the bone marrow and nerves can cause anemia, infections, vision loss, or other problems.

\medterm{Osteophyte} A bony outgrowth forming at the edge of a joint, commonly with osteoarthritis. It may cause stiffness, pain, reduced movement, or pressure on nearby nerves, although some osteophytes cause no symptoms.

\medterm{Osteoplasty} Surgery to repair, reshape, or reconstruct bone to restore structure, movement, or function. The technique depends on the bone problem and may be combined with fixation, grafting, or joint surgery.

\medterm{Osteopoikilosis} Osteopoikilosis is an inherited benign bone condition with multiple small dense spots on imaging, usually asymptomatic and discovered incidentally.

\medterm{Osteopontin} A protein involved in bone remodeling, immune responses, inflammation, and attachment between cells and tissues. Its level may change in several diseases, but a measurement is not by itself a specific diagnosis.
\medterm{Osteoporosis} A bone disease in which bone is broken down faster than it is replaced, or its structure and strength change. The bones become less dense and more likely to break, sometimes after a minor fall or other low-impact force. Osteoporosis often causes no symptoms until a fracture occurs. Fractures commonly affect the hip, spine, and wrist; weakened bones in the spine may collapse and cause back pain or loss of height. Osteoporosis can affect anyone, but it is more common in females, especially after menopause, when lower estrogen levels can speed bone loss. It can also affect males. Risk is influenced by age, family history, hormones, inadequate calcium or vitamin D, some medicines, smoking, alcohol, and other health conditions.

\medterm{Osteoporosis Fracture Risk Assessment} An evaluation of future fracture risk using bone density, age, previous fractures, falls, medicines, family history, and other health factors. The result helps determine whether lifestyle measures, monitoring, or bone-strengthening treatment is appropriate.

\medterm{Osteoporotic Fracture} A fracture resulting from weakened, porous bone after a low-energy event that would not usually break healthy bone.

\medterm{Osteoradionecrosis} Death of bone within an area treated with radiation, most often the jaw after head-and-neck cancer treatment. The bone may become exposed or fail to heal after injury or dental work and can become infected.
\medterm{Osteosarcoma} An aggressive cancer that begins in bone-forming cells and produces abnormal immature bone. It often develops near the ends of long bones in children and young adults and may cause persistent pain, swelling, or a fracture.

\medterm{Osteosarcoma, Secondary} An osteogenic sarcoma arising in previously abnormal or treated bone, such as bone affected by Paget disease, prior radiation, or chronic bone lesions; prognosis and treatment depend on stage and resectability.
\medterm{Osteosclerosis} Abnormally increased bone density seen on imaging or examination. It may result from inherited disease, altered bone renewal, chronic stress, inflammation, medicines, or cancer; dense bone can still be weak or painful.

\medterm{Osteosynthesis} Surgical repair and stabilization of a broken bone using implants such as plates, screws, rods, nails, or wires. The implants hold the pieces in alignment while the bone heals.

\medterm{Osteotome} A sharp surgical instrument used to cut or shape bone.

\medterm{Osteotomy} A surgical procedure involving cutting and repositioning of bone to correct deformity,
realign joint surfaces, alter weight-bearing forces, or change limb length.

\medterm{Osteotomy Of The Knee} Surgery that cuts and reshapes a knee-region bone to correct alignment and redistribute weight across the joint.

\medterm{Ostium} A natural opening into or out of an organ, cavity, blood vessel, or duct. It allows air, fluid, or other contents to pass between connected body spaces.
\medterm{Ostomy} A surgically created opening that connects part of the bowel or urinary tract to the skin so stool or urine can leave the body. A colostomy or ileostomy passes stool; a urostomy passes urine. A pouch usually collects the output and protects the surrounding skin.
\medterm{Ostomy Appliance} A pouching system that collects stool or urine from a surgically created opening in the abdomen. Colostomy and ileostomy appliances collect stool; urostomy appliances collect urine. The skin around the opening needs regular protection and monitoring.

\medterm{Ostraceous Psoriasis} A rare severe form of psoriasis in which plaques become very thick, hard, and layered, resembling an oyster shell. It can be difficult to treat and may require specialist care and medicines affecting the whole immune system.

\medterm{Otalgia} Pain in or around the ear. It may come from an ear infection, wax, injury, or pressure problem, or may be referred from the teeth, jaw, throat, or neck. Persistent pain, discharge, hearing loss, fever, or severe pain needs medical assessment.
\medterm{Otalgia (Earache)} Pain in or around the ear, arising from the ear itself or referred from structures such as the teeth, jaw, throat, or cervical spine.

\synonyms
Earache

\medterm{Othello Syndrome} A delusional belief that a partner is unfaithful despite insufficient evidence; it may occur with psychiatric, neurologic, or substance-related disorders.


\medterm{Other Pain Conditions} Pain conditions that do not fit neatly into a single category, including nociplastic pain, fibromyalgia, chronic
tension-type headache, temporomandibular pain, chronic pelvic pain, and complex regional
pain syndrome. Diagnosis and treatment depend on the specific condition.

\medterm{Other Race} Other race is a demographic category for a person who does not select or fit the available listed race categories; classification varies by system and should not be treated as a biological diagnosis.

\medterm{Otic Barotrauma} Injury to the ear caused by failure to equalize pressure during flying, diving, or other pressure changes.
\medterm{Otic Capsule} The dense part of the skull bone that surrounds and protects the inner ear, including the hearing and balance organs. Disease or abnormal remodeling in this bone can affect hearing or balance.

\medterm{Otic Ointment} A semisolid medicine formulated for application to the outer ear or ear canal; it should be used only as directed, especially when the eardrum may be perforated.

\medterm{Otitis} Inflammation of the ear, which may be infectious or noninfectious. Otitis externa affects the ear canal, otitis media affects the middle ear, and inner-ear inflammation may cause vertigo or hearing loss. Symptoms can include pain, fever, blockage, discharge, or reduced hearing and vary with the location and cause.
\medterm{Otitis (Ear Infection)} Alternate index label for Otitis.

\medterm{Otitis Externa} Inflammation or infection of the outer ear canal, commonly causing ear pain, itching, tenderness, swelling, or discharge. Keeping water and objects out of the ear can help while it heals.
\synonyms
Swimmers Ear (Otitis Externa)

\medterm{Otitis Media} Inflammation or infection in the middle ear, often after a cold or other upper-airway infection. It may cause ear pain, fever, reduced hearing, irritability, or fluid behind the eardrum.
\medterm{Otitis Media With Effusion} Fluid behind the eardrum without the severe pain, fever, or pus of an acute ear infection. It may cause blocked hearing or speech and learning problems in children and often clears, but persistent fluid needs hearing assessment.

\medterm{Oto-Rhino-Laryngology} Another name for otolaryngology, the specialty of the ear, nose, throat, head, and neck.
\medterm{Otoacoustic Emissions} Very soft sounds naturally produced by a healthy inner ear that can be measured with a small probe placed in the ear canal. The test is quick and painless and is commonly used to screen newborn hearing and assess inner-ear function.

\synonyms
OAE; otoacoustic-emission testing; cochlear emissions.

\medterm{Otocephaly} A severe congenital craniofacial malformation characterized by absent or markedly
underdeveloped mandible, ventral displacement or fusion of the ears, and associated
facial and airway abnormalities. It is usually lethal because of profound impairment of
breathing and feeding.

\medterm{Otolaryngologist} An otolaryngologist is a physician specializing in disorders of the ear, nose, throat, head, neck, voice, swallowing, and related surgery.

\medterm{Otolaryngology} The medical specialty treating disorders of the ears, nose, throat, head, and neck. Specialists manage hearing, balance, breathing, swallowing, voice, sinus, facial, and related surgical problems.
\medterm{Otolith} A tiny calcium-containing crystal in the inner ear that helps sense gravity and head movement. If it moves into the wrong balance canal, it can cause brief spinning attacks known as positional vertigo.

\medterm{Otoliths} Tiny calcium-carbonate crystals in inner-ear organs that sense gravity and straight-line movement. If they move into the wrong balance canal, they can cause brief spinning attacks called positional vertigo.
\medterm{Otology} The study and medical care of the ear and hearing.
\medterm{Otoplasty} Surgery to alter the shape, position, or proportion of the external ear, commonly to correct prominent or congenital ear differences.

\medterm{Otorhinolaryngology} The medical specialty concerned with diseases of the ear, nose, throat, head, and neck; also called otolaryngology.

\synonyms
otolaryngology.

\medterm{Otorrhea} Fluid, pus, blood, or other discharge flowing from the ear. Causes include infection, eardrum perforation, injury, or a growth; persistent, bloody, foul-smelling, or clear discharge needs medical assessment.

\medterm{Otorrhoea} Fluid, pus, or blood draining from the ear. Common causes include an ear-canal infection, a middle-ear infection with a perforated eardrum, injury, or a growth. Clear watery discharge after a head injury is an emergency because it may be brain fluid and can lead to meningitis.
\medterm{Otosclerosis} An abnormal hardening and fixation of middle-ear bones that can gradually reduce hearing by blocking sound transmission.
\medterm{Otoscope} A lighted instrument used to examine the ear canal and eardrum. It can help identify earwax, infection, a foreign object, fluid, swelling, or a perforated eardrum. Examination findings are interpreted with symptoms and, when needed, hearing tests.
\medterm{Otoscopy} Looking into the ear canal and at the eardrum with a lighted instrument called an otoscope. It can help find earwax, infection, a foreign object, a hole in the eardrum, fluid, or other causes of ear symptoms.
\medterm{Ototoxicity} Damage to the inner ear from a medicine, chemical, or other toxic exposure. It may cause hearing loss, ringing in the ears, dizziness, or imbalance and can be temporary or permanent.

\synonyms
Auditory toxicity; vestibular toxicity; ototoxic effect.

\medterm{Ounce} An ounce is a unit of mass or fluid volume whose exact conversion depends on the measurement system; medication instructions should use unambiguous units.

\medterm{Outcome} The result or final effect of a disease, treatment, exposure, or healthcare intervention. Outcomes may include recovery, symptom control, complications, disability, quality of life, or death and can be measured in different ways.
\medterm{Outcome Measurement} The use of defined measures to assess a patient’s symptoms, function, quality of life, or other results before and after care.



\medterm{Outer Ear} The visible auricle and the external auditory canal, which collect and conduct sound toward the eardrum.

\medterm{Outer Ear Infection} Alternate index label; see \textit{Swimmer's ear}.

\medterm{Outpatient} A person who receives medical evaluation or treatment without being admitted overnight to a hospital or other inpatient facility.
\medterm{Outpatient Rehabilitation} Rehabilitation provided without hospital admission to improve function after illness, injury,
surgery, or disability. The setting, frequency, and duration are individualized.

\medterm{Output} The amount of a substance produced or eliminated by the body or by an organ over a specified time, such as urine output or cardiac output.
\medterm{Ova} Plural of ovum; eggs or egg cells, and in parasitology, eggs of worms or other parasites.

\medterm{Ovalocyte Count} An ovalocyte count estimates the number or proportion of oval-shaped red blood cells on a blood smear and may help evaluate anemia.

\medterm{Ovarian} Relating to an ovary. The ovaries release eggs and produce hormones. Ovarian problems include cysts, inflammation, reduced function, twisting that cuts off blood supply, and tumors; sudden severe pelvic pain may indicate twisting and needs emergency care.
\medterm{Ovarian Cancer} Cancer arising in an ovary or nearby reproductive tissue. It may cause persistent bloating, pelvic or abdominal pain, early fullness, urinary changes, or abnormal bleeding, but early disease may have few symptoms.

\synonyms
Cancer, Ovarian

\medterm{Ovarian Cyst} Alternate index label for Ovarian Cysts.

\medterm{Ovarian Cysts} Fluid-filled or partly solid sacs in or on an ovary. Many are related to the menstrual cycle and resolve spontaneously, while others may be associated with endometriosis, benign tumors, or hormonal conditions. Ovarian cysts may cause pelvic pain, menstrual changes, pressure, bleeding, rupture, or twisting of the ovary; some cause no symptoms.

\synonyms
Cyst, Ovarian
Cysts Ovary (Ovarian Cysts)

\medterm{Ovarian Failure} Reduced ovarian activity causing irregular or absent periods, difficulty becoming pregnant, and symptoms of low estrogen such as hot flushes or vaginal dryness. When it occurs before age 40, it is usually called primary ovarian insufficiency and needs evaluation.

\medterm{Ovarian Fibromata} Usually noncancerous tumors made of fibrous tissue that develop in an ovary. They may cause no symptoms or produce pelvic pressure, pain, menstrual changes, or fluid in the abdomen when large.

\medterm{Ovarian Follicle} A small fluid-filled structure in an ovary that contains an immature egg and supporting cells. During each menstrual cycle, one follicle usually matures and releases an egg during ovulation.

\medterm{Ovarian Gonadoblastoma} A rare tumor containing cells that normally develop into eggs and cells that support sex-hormone production. It often occurs in an atypically developed gonad, may contain Y-chromosome material, and can produce hormones.

\medterm{Ovarian Hemorrhage} Bleeding into or around an ovary, often from a ruptured cyst or blood vessel. It may cause sudden pelvic pain, weakness, dizziness, or internal blood loss and can require urgent assessment or surgery.

\medterm{Ovarian Hyperstimulation Syndrome} A complication of fertility treatment in which the ovaries enlarge and fluid moves into the abdomen or other tissues. It may cause bloating, pain, vomiting, reduced urine, breathlessness, blood clots, or kidney problems and can be severe.

\synonyms
OHSS

\medterm{Ovarian Infection} Infection involving an ovary, often as part of pelvic inflammatory disease or an abscess, causing pelvic pain, fever, and abnormal discharge.

\medterm{Ovarian Metastasis} Cancer that has spread to an ovary from another organ, commonly the digestive tract, breast, or reproductive organs. It may resemble a primary ovarian cancer, so tissue testing is important because treatment depends on the original site.

\medterm{Ovarian Overproduction Of Androgens} Ovarian androgen overproduction is excessive production of male-pattern hormones by the ovaries, which may cause acne, excess facial hair, irregular periods, or infertility.

\medterm{Ovarian Pregnancy} A rare ectopic pregnancy that implants in an ovary rather than inside the uterus or fallopian tube. It may cause one-sided pelvic pain or internal bleeding and usually needs urgent ultrasound assessment and treatment.

\medterm{Ovarian Rupture} Ovarian rupture is tearing of an ovary or ovarian cyst, which may cause sudden pelvic pain and internal bleeding and sometimes requires urgent treatment.

\medterm{Ovarian Teratoma} A tumor of an ovary made from reproductive cells that can form different kinds of tissue, such as skin, hair, or bone. Mature teratomas are usually noncancerous; immature teratomas can be cancerous.

\medterm{Ovarian Thecoma} An ovarian thecoma is a usually benign sex-cord stromal tumor that may produce estrogen and cause abnormal uterine bleeding or other hormonal effects.

\medterm{Ovarian Torsion} Twisting of an ovary, often together with its fallopian tube, that can cut off blood flow. It causes sudden one-sided pelvic pain, often with nausea or vomiting, and is a gynecologic emergency requiring urgent assessment and usually surgery.

\medterm{Ovariectomy (Oophorectomy)} Surgery to remove one or both ovaries. It may treat ovarian disease or reduce cancer risk. Removing both ovaries before natural menopause causes permanent loss of ovarian hormone production.
\medterm{Ovaries} The paired female gonads that produce oocytes and secrete sex hormones, mainly estrogen and progesterone. They are located on either side of the uterus and participate in the menstrual cycle and reproduction.

\synonyms
ovary

\medterm{Ovariotomy} An older term for surgical removal or incision of an ovary.
\medterm{Ovary} A reproductive organ that releases eggs and produces hormones such as estrogen and progesterone. The ovaries help control menstrual cycles, ovulation, pregnancy preparation, and many body changes during reproductive life.
\medterm{Ovary Disease} Any disorder affecting an ovary, including cysts, infection, inflammation, tumors, hormonal dysfunction, torsion, or bleeding.
\medterm{Ovary Neoplasm} An abnormal growth in an ovary that may be noncancerous, have uncertain behavior, or be cancerous. Symptoms can include bloating, pelvic pain, pressure, or abnormal bleeding, but some growths are found incidentally on imaging.

\medterm{Oven Cleaner Poisoning} Oven cleaners can severely burn the mouth, throat, stomach, skin, or eyes because they may contain strong alkalis. Do not induce vomiting; rinse exposed skin or eyes with water and contact emergency services or poison control.

\medterm{Over-Hydration} Excess accumulation of body water, which may cause edema, dilutional hyponatremia, or heart and lung complications.
\medterm{Over-Nutrition} Excessive intake or accumulation of energy or nutrients that can contribute to obesity, metabolic disease, or toxicity.
\medterm{Over-The-Counter} A medicine available without a prescription for self-treatment of specified common problems. It can still cause harm, interact with prescribed medicines, or be unsafe with certain illnesses, so the label and dose should be followed carefully.

\medterm{Over-The-Counter Medications} Medicines sold without a prescription for selected common symptoms or conditions. They are not risk-free: incorrect use, prolonged use, interactions, allergies, or underlying illness can cause serious harm.

\medterm{Over-The-Counter Medicines} Over-the-counter medicines can be bought without a prescription for common problems, but they may still cause side effects or interact with prescribed medicines.

\medterm{Over-The-Counter Pain Relievers} Over-the-counter pain relievers include medicines such as acetaminophen and nonsteroidal anti-inflammatory drugs; choice depends on age, health conditions, and other medicines.

\medterm{Overactive Bladder} A pattern of a sudden, difficult-to-delay need to urinate, usually with frequent urination and waking at night, with or without leakage. Causes include bladder irritation, nerve problems, obstruction, medicines, and sometimes no identifiable cause.

\synonyms
OAB
OAB (Overactive Bladder)

\medterm{Overall Survival} The length of time from diagnosis, treatment, or study entry until death from any cause. Doctors and researchers use it to describe how long people live and to compare treatments. It differs from time without worsening disease because it includes survival regardless of whether the illness has progressed.

\medterm{Overbite} Vertical overlap of the upper front teeth over the lower front teeth; excessive overbite can damage gums, teeth, or jaw function.

\medterm{Overcoming Breastfeeding Problems} Breastfeeding problems such as pain, poor latch, low supply, engorgement, or blocked ducts can often improve with feeding assessment, positioning help, and timely lactation or medical support.

\medterm{Overdenture} A removable denture supported by dental implants or remaining natural teeth. This support can improve stability, chewing, speech, and comfort, while attachments may help preserve jawbone and reduce denture movement.

\medterm{Overdose} An overdose is exposure to a medicine, drug, or substance in an amount capable of causing toxicity, with effects determined by dose, formulation, route, timing, and patient factors.
\medterm{Overflow Urinary Incontinence} Urine leakage caused by a bladder that remains too full because it does not empty properly. Weak flow, difficulty starting, frequent small amounts, dribbling, or constant leakage may occur; blockage or nerve disease can be responsible.

\medterm{Overgrowth Syndrome} A group of conditions present from birth in which the whole body, one side, or particular tissues grow more than expected. Some forms also cause developmental differences or increase the risk of certain tumors.

\medterm{Overlap Myositis} Muscle inflammation occurring together with another autoimmune connective-tissue disease, such as lupus, systemic sclerosis, mixed connective-tissue disease, or Sjögren syndrome. It may cause muscle weakness, skin changes, joint symptoms, swallowing problems, or lung disease; treatment is tailored to the affected organs.

\medterm{Overlay Denture} A removable denture that fits over remaining teeth, roots, or implants to improve support, chewing, and retention.

\medterm{Overload Iron (Iron Overload)} Alternate index label for Iron Overload.

\medterm{Overriding Aorta} A birth-related heart abnormality in which the main artery leaving the heart sits over a hole between the lower chambers and receives blood from both. It commonly occurs as part of tetralogy of Fallot.
\medterm{Oversedate} To give excessive sedation, causing unusually deep sleep, slowed responses, impaired breathing, or difficulty maintaining the airway.

\medterm{Overt Proteinuria} Abnormally high protein loss in the urine that is detectable by routine testing and may indicate kidney disease.

\medterm{Overuse Syndrome (Repetitive Motion Disorders (RMDs))} Alternate index label for Repetitive Motion Disorders (RMDs).

\medterm{Overweight} Body weight above the usual range for height. In adults, a Body Mass Index of 25 to less than 30 kg/m\(^2\) is commonly classified as overweight, but it does not directly measure body fat.
\medterm{Oviduct} A tube carrying an ovum from the ovary toward the uterus; in humans it is the uterine or fallopian tube.
\medterm{Ovomucin} A glycoprotein in egg white that contributes to its gel-like viscosity and helps stabilize the egg-white structure.

\medterm{Ovotestis} A gonad containing both ovarian-type and testicular tissue. It is an uncommon difference of sex development and may produce ovarian hormones, testicular hormones, eggs, sperm, or none of these; evaluation is individualized.
\medterm{Ovulation} Release of a mature egg from an ovarian follicle, usually around the middle of the menstrual cycle.
\medterm{Ovulation Home Test} An ovulation home test detects the rise in urinary luteinizing hormone that usually occurs shortly before ovulation.

\medterm{Ovum} The female reproductive cell released from an ovary during ovulation. If a sperm enters it, the cell completes fertilization and joins with the sperm’s genetic material to form the first cell of an embryo.
\medterm{Oxacillin} A penicillinase-resistant penicillin used mainly to test for methicillin resistance in staphylococci; it is not commonly used as the preferred treatment for most infections.

\medterm{Oxalate} A salt or ester of oxalic acid that can form kidney stones and binds calcium.
\medterm{Oxalates} Oxalates are salts or esters of oxalic acid that can bind calcium and form calcium-oxalate kidney stones; excess dietary absorption or inherited hyperoxaluria can cause renal injury.

\medterm{Oxalic Acid} A dicarboxylic acid whose oxalate salts can contribute to calcium-oxalate kidney stones; concentrated exposure is toxic.
\medterm{Oxalic Acid Poisoning} Swallowing oxalic acid can burn the digestive tract and bind calcium, causing muscle cramps, abnormal heart rhythms, seizures, kidney injury, or shock. It is a medical emergency requiring immediate poison-control or emergency care.

\medterm{Oxaliplatin} Oxaliplatin is a platinum chemotherapy drug that forms DNA cross-links and is used especially in colorectal cancer regimens; acute cold-triggered neuropathy and cumulative sensory neuropathy are characteristic toxicities.

\medterm{Oxaloacetate} A small molecule used in the cell’s energy-making cycle. It combines with another molecule to begin a step that produces energy and can also help make glucose and the amino acid aspartate.

\medterm{Oxandrolone} An anabolic steroid used in selected conditions involving severe weight loss or catabolic stress; it can affect the liver, lipids, and hormone function.

\medterm{Oxaprozin} A nonsteroidal anti-inflammatory drug used to relieve pain and inflammation in conditions such as osteoarthritis and rheumatoid arthritis.

\medterm{Oxazepam} A benzodiazepine medicine used short-term for anxiety or alcohol withdrawal, causing sleepiness and possible dependence.
\medterm{Oxazepam Overdose} Taking too much oxazepam, a sedative medicine, can cause severe sleepiness, confusion, poor coordination, slowed breathing, coma, or death, especially with alcohol or other sedatives. Immediate emergency advice is needed.

\medterm{Oxazole} Oxazole is a five-membered chemical ring containing oxygen and nitrogen that is used as a building block in medicinal and other chemistry.

\medterm{Oxcarbazepine} Oxcarbazepine is an antiseizure medicine that blocks voltage-gated sodium channels and treats focal seizures; dizziness, ataxia, rash, and clinically important hyponatraemia may occur.

\medterm{Oxidant} A substance that can damage cells by taking electrons from other molecules. Oxidants are produced normally during metabolism and inflammation, but excessive amounts can cause oxidative stress and tissue injury.

\synonyms
Oxidizing agent

\medterm{Oxidants} Substances or reactive molecules that accept electrons or promote oxidation; excessive oxidant activity can damage proteins, lipids, and DNA.

\medterm{Oxidation} A chemical process that changes a substance through loss of electrons or reaction with oxygen.
\medterm{Oxidation-Reduction} A chemical reaction involving transfer of electrons, in which one substance is oxidized and another is reduced.

\medterm{Oxidative Phosphorylation} The final stage of energy production in mitochondria, where electrons from food drive the movement of hydrogen ions and power the enzyme that makes ATP, the main energy-carrying molecule in cells.

\medterm{Oxidative Stress} A harmful imbalance between reactive oxygen-containing molecules and the body’s protective antioxidants. Excessive oxidative stress can injure cell membranes, proteins, and DNA and may contribute to inflammation, aging, or disease, although the term alone does not identify a specific illness.
\medterm{Oxidoreductase} An enzyme that transfers electrons or hydrogen between molecules during a chemical reaction. These enzymes are important in energy production, drug breakdown, and many other cell processes.

\medterm{Oxime} A type of chemical compound. Some oximes, such as pralidoxime, are used as emergency antidotes because they can restore an important nerve enzyme after organophosphate pesticide poisoning.

\medterm{Oximeter} A device that estimates blood oxygen level, usually by shining light through a finger, toe, or earlobe. A pulse oximeter also displays pulse rate, but readings can be affected by movement, poor circulation, or skin factors.
\medterm{Oximetry} Measurement of oxygen levels in blood or tissue. A pulse oximeter estimates blood oxygen through a skin sensor, while arterial testing measures oxygen more directly; results can be affected by poor circulation or motion.


\medterm{Oxter} Oxter is an anatomical term for the armpit or axilla.

\medterm{Oxybuprocaine} A local anaesthetic eye drop that briefly numbs the surface of the eye. It may be used for examinations or short procedures and should not be used repeatedly without supervision because it can delay healing and damage the cornea.
\medterm{Oxycephaly} A cone-shaped skull caused by premature fusion of cranial sutures, also called acrocephaly.
\medterm{Oxycodone} Oxycodone is an opioid analgesic used for moderate-to-severe pain; sedation, constipation, respiratory depression, tolerance, dependence, and overdose risk increase with alcohol or other sedatives.

\medterm{Oxycodone Hydrochloride} The hydrochloride salt of oxycodone, an opioid analgesic used for moderate-to-severe pain; it can cause sedation, constipation, dependence, and fatal respiratory depression in overdose.

\medterm{Oxygen} A gas essential for aerobic metabolism and tissue energy production.
\medterm{Oxygen Consumption} The amount of oxygen the body uses in a given time. It reflects the body’s energy needs and rises with exercise, fever, or illness; it can be measured during exercise testing to assess heart and lung function.

\medterm{Oxygen Content} The amount of oxygen carried in blood, mainly bound to hemoglobin with a small dissolved component.
\medterm{Oxygen Delivery} The amount of oxygen transported from the lungs to body tissues each minute. It depends mainly on cardiac output, the amount of hemoglobin, and how much oxygen the hemoglobin carries. Inadequate delivery can injure organs when the body cannot compensate.

\medterm{Oxygen Desaturation Index} The average number of oxygen desaturation events per hour of sleep, usually reported from a sleep study as the ODI; interpretation depends on the recording method and clinical context.

\medterm{Oxygen Enhancement Ratio} A measure of how much oxygen increases the effect of radiation on cells. It compares the radiation dose needed with little oxygen to the dose needed when oxygen is normal to produce the same amount of cell injury.
\medterm{Oxygen Mask} A mask that delivers supplemental oxygen through the nose and mouth. The type and flow rate are chosen according to oxygen needs; oxygen should be used as prescribed and monitored.

\medterm{Oxygen Pulse} The amount of oxygen used by the body with each heartbeat, estimated during exercise as oxygen uptake divided by heart rate.

\medterm{Oxygen Safety} Practices that reduce fire, explosion, pressure, and treatment risks during oxygen use, including keeping oxygen away from flames, smoking, oils, and unsafe equipment.

\medterm{Oxygen Saturation} The percentage of hemoglobin in the blood carrying oxygen. A finger or ear sensor estimates it without taking blood, while an arterial blood test measures it more directly; the healthy target varies with the person’s condition.

\medterm{Oxygen Tent} A canopy or enclosure used to deliver oxygen around a patient’s head; modern oxygen masks, nasal cannulas, or ventilatory devices are usually more precise.

\medterm{Oxygen Therapy} Giving extra oxygen through a tube, mask, or breathing machine when the blood oxygen level is too low or is expected to fall. The amount is prescribed and monitored because too much oxygen can also be harmful; oxygen must be kept away from flames.
\medterm{Oxygen Therapy In Infants} Carefully measured extra oxygen given to a baby whose blood oxygen is too low or whose lungs are not exchanging gases adequately. Clinicians monitor oxygen closely because both too little and too much can injure developing organs.

\medterm{Oxygen Toxicity} Tissue injury caused by exposure to oxygen at excessively high partial pressures or concentrations, potentially affecting the lungs, eyes, or central nervous system.

\medterm{Oxygen Transport} The movement of oxygen from the lungs to tissues, mainly bound to hemoglobin and partly dissolved in plasma, with delivery affected by blood flow and hemoglobin concentration.

\medterm{Oxygen-Hemoglobin Dissociation Curve} A graph showing how strongly hemoglobin holds oxygen at different oxygen levels. Temperature, acidity, and carbon dioxide change this binding and therefore affect how readily oxygen is released to tissues.

\medterm{Oxygenases} Enzymes that incorporate oxygen from molecular oxygen into a substrate; monooxygenases add one oxygen atom and dioxygenases add both.

\medterm{Oxygenated Blood} Blood that has picked up oxygen in the lungs and carries it through arteries to body tissues. It is usually bright red, while blood returning to the lungs contains less oxygen and more carbon dioxide.

\medterm{Oxygenation} The process of supplying oxygen to blood and tissues, or the degree to which hemoglobin and tissues are adequately oxygenated.

\medterm{Oxygenator} A machine that adds oxygen to blood and removes carbon dioxide when the lungs cannot do this adequately. It is used during heart-lung bypass surgery and some forms of extracorporeal life support.

\medterm{Oxyhaemoglobin} Hemoglobin with oxygen attached to it. It forms mainly in the lungs and carries oxygen through the bloodstream to body tissues.
\medterm{Oxyhemoglobin} Hemoglobin bound to oxygen, formed mainly in the lungs and carrying oxygen through arterial blood to tissues.

\medterm{Oxylipins} Oxylipins are oxygenated fatty-acid derivatives that act as local mediators of inflammation, vascular tone, pain, and other physiologic responses.

\medterm{Oxymetazoline} A topical alpha-adrenergic agonist used as a nasal decongestant or ocular vasoconstrictor; prolonged nasal use can cause rebound congestion.
\medterm{Oxymetholone} Oxymetholone is an anabolic steroid used in selected severe anemias; it can cause liver injury, fluid retention, and hormonal side effects.

\medterm{Oxymorphone} A potent opioid analgesic used for severe pain; it can cause sedation, constipation, dependence, and life-threatening respiratory depression.

\medterm{Oxyphil Cell} A cell with many mitochondria, giving it a distinctive granular appearance under a microscope. Oxyphil cells occur in tissues such as the parathyroid and thyroid, where their significance depends on the organ and the surrounding findings.

\medterm{Oxyphil Cell Carcinoma} Alternate index label; see \textit{Hurthle cell cancer}.

\medterm{Oxyphilic Adenoma} An oxyphilic adenoma is a usually benign glandular tumor composed of cells rich in mitochondria, such as an oncocytic parathyroid or salivary-gland tumor.

\medterm{Oxytocic} Causing the womb to contract. Oxytocic medicines, such as oxytocin, may start or strengthen labor or help control heavy bleeding after birth; they require monitoring because overly strong or frequent contractions can harm the pregnant person or baby.
\medterm{Oxytocics} Medicines that stimulate uterine contractions, used to start or strengthen labor and to prevent or treat heavy bleeding after birth. They require monitoring because overly strong or frequent contractions can harm the pregnant person or baby.

\medterm{Oxytocin} A hormone made in the brain and released by the pituitary gland. It helps the uterus contract during labor and causes milk to flow during breastfeeding; a manufactured form is used under monitoring to induce labor or control bleeding after birth.
\medterm{Oxytocin Infusion} A controlled drip of synthetic oxytocin into a vein to start or strengthen labor or help control heavy bleeding after birth. It requires continuous assessment because overly strong or frequent contractions can harm the pregnant person or baby.

\medterm{Oxyuriasis} Infection with pinworms, especially Enterobius vermicularis. The worms commonly cause itching around the anus at night, disturbed sleep, and sometimes abdominal discomfort; treatment usually includes medicine and household hygiene.

\medterm{Ozaena} A chronic form of atrophic rhinitis in which the inside of the nose becomes thin and widened, with dry crusts, a foul smell, and a blocked sensation despite an open passage. Medical assessment is needed to exclude infection and other causes.

\medterm{Ozena} Ozena is a chronic form of atrophic rhinitis with widened nasal passages, crusting, foul odor, and impaired mucosal function.

\medterm{Ocular Larva Migrans} Ocular toxocariasis caused by migrating larvae of \textit{Toxocara} roundworms, usually affecting one eye. It can cause retinal or intraocular inflammation, a white pupil, reduced vision, strabismus, or retinal damage and requires ophthalmic assessment; diagnosis is based on the eye findings, exposure history, and supportive serology.

\medterm{Omega-3 Fatty Acid} A type of polyunsaturated fat found in fish, seafood, flaxseed, chia, walnuts, and some oils. Omega-3s support cell membranes and cardiovascular function; supplements can interact with medicines and are not a substitute for treatment.

\medterm{Opioids} Medicines and drugs that act on opioid receptors to relieve pain. They can cause sleepiness, constipation, dependence, and life-threatening breathing suppression, especially with alcohol or sedatives; overdose may be reversed with naloxone.

\medterm{Opsoclonus-Myoclonus Syndrome} A rare neurologic disorder causing rapid, involuntary eye movements, muscle jerks, poor coordination, and sometimes behavioral or cognitive changes. In children it may be associated with neuroblastoma or an immune reaction.

\medterm{Optic Disc Swelling} Enlargement or elevation of the optic nerve head where it enters the eye. Causes range from raised pressure inside the skull to inflammation, ischemia, or local eye disease; sudden vision changes require urgent assessment.

\medterm{Oral-Facial-Digital Syndrome} A group of genetic disorders affecting the mouth, face, fingers, and toes. Features vary and may include oral abnormalities, unusual facial development, shortened or joined digits, and neurologic or kidney problems.

\medterm{Orbital Blowout Fracture} A break in one of the thin bones forming the eye socket, usually after blunt trauma. It can cause swelling, double vision, numbness, or restricted eye movement and may need urgent evaluation.

\medterm{Orthodontic} Relating to the diagnosis, prevention, or treatment of problems with tooth and jaw position. Orthodontic treatment may use braces, aligners, retainers, or other appliances.

\medterm{Orthorexia Nervosa} A term used for an unhealthy preoccupation with eating foods considered pure or healthy, leading to severe restriction, distress, or impaired daily life. It is not yet a formal standalone diagnosis in major classification systems, and professional assessment is important.

\medterm{Ozone} Ozone is a reactive form of oxygen that can irritate and injure the airways and lungs when inhaled, despite useful atmospheric and disinfecting roles.
\medterm{Oral Route} The administration of a medicine through the mouth for swallowing and absorption through the gastrointestinal tract. Clear instructions should spell out “by mouth” to reduce medication-route errors.

\synonyms
P.O.; per os; by mouth
